\SourcePage{221}
\SGAsection{5}{诱导拓扑斯}
\textbf{5.1。} 设 $E$ 为一个\Term{topos}{拓扑斯}，$X$ 为 $E$ 的一个\Term{object}{对象}。则 $E$ 中位于 $X$ 之上的\Term{object}{对象}的\Term{category}{范畴} $E_{/X}$ 是一个\Term{topos}{拓扑斯}，例如这立即由 Giraud 判据 1.2 ii) 得出。也可以利用 1.2.1，把 $E$ 实现为\Term{sheaf}{层}\Term{category}{范畴} $C^\sim$，
\SourcePage{222}
其中 $C$ 是 $E$ 的一个\Term{generating-subcategory}{生成子范畴}，并且可以选择使 $X\in\operatorname{ob}C$；此时我们知道 $E_{/X}=C^\sim_{/X}$ \Term{equivalence}{等价}于 $(C_{/X})^\sim$（第 III 讲 5.4），因而是\Term{topos}{拓扑斯}。

称\Term{topos}{拓扑斯} $E_{/X}$ 为在 $E$ 的\Term{object}{对象} $X$ 上的\Term{induced-topos}{诱导拓扑斯}。

\textbf{5.2。} 我们要定义一个\Term{canonical}{典范}的\Term{topos-morphism}{拓扑斯态射}
\[j_X:E_{/X}\longrightarrow E,\tag{5.2.1}\]
称为\Term{induced-topos}{诱导拓扑斯} $E_{/X}$ 到周围\Term{topos}{拓扑斯} $E$ 的包含\Term{morphism}{态射}，或者更恰当地（见 5.7），称为 $E$ 在 $X$ 处的\Term{localization-morphism}{局部化态射}。这个\Term{morphism}{态射}\Term{correspondence}{对应}于三个依次\Term{adjunction}{伴随}的\Term{functor}{函子}（见 3.1.3）
\[j_{X!},\quad j_X^*,\quad j_{X*},\tag{5.2.2}\]
它们可具体说明如下。

 a) \Term{functor}{函子} $j_{X!}:E_{/X}\to E$ 是从 $E_{/X}$ 到 $E$ 的“遗忘结构\Term{arrow}{箭头}”\Term{functor}{函子}。

 b) \Term{functor}{函子} $j_X^*:E\to E_{/X}$ 定义为
\[j_X^*(Z)=(X\times Z,\operatorname{pr}_1),\]
其中 $\operatorname{pr}_1:X\times Z\to X$ 是第一\Term{projection}{投影}。也可以把它解释为关于\Term{morphism}{态射} $X\to e_E$ 的\Term{base-change-functor}{换基函子}，其中 $e_E$ 为 $E$ 的\Term{final-object}{终对象}。由\Term{base-change-functor}{换基函子}的定义，$j_X^*$ 显然是 $j_{X!}$ 的\Term{right-adjoint}{右伴随}。特别地，它与\Term{projective-limit}{投射极限}交换。由\Term{topos}{拓扑斯}所特有的\Term{exactness}{正合性}（第 II 讲 4），它也与\Term{inductive-limit}{归纳极限}交换。

 c) 由后一个事实可知（1.6），$j_X^*$ 有一个\Term{left-adjoint}{左伴随}
\[j_{X*}:E_{/X}\longrightarrow E.\]
对于 $X$ 上的\Term{object}{对象} $X'$，\Term{object}{对象} $j_{X*}(X')$ 有时也用下列某个记号表示：
\[\prod_{X/e_E}(X'/X),\qquad \mathscr{H}om_{X/e_E}(X,X'),\qquad \operatorname{Res}_{X/e_E}(X'/X)\]
（“Weil \Term{restriction}{限制}”）。本书博学的读者无疑已经熟悉其中的某一个。

\SourcePage{223}
\textbf{5.3。} 应当注意，\Term{functor}{函子} $j_{X!}$ 与\Term{fiber-product}{纤维积}交换，并把\Term{monomorphism}{单态射}变为\Term{monomorphism}{单态射}，但一般并不因此而\Term{left-exact}{左正合}。它只有在 $X$ 为 $E$ 的\Term{final-object}{终对象}时才\Term{left-exact}{左正合}（因为 $X=j_{X!}(X,\operatorname{id}_X)$，且 $(X,\operatorname{id}_X)$ 为 $E_{/X}$ 的\Term{final-object}{终对象}），也就是说，当且仅当 $j_X$ 实际上是\Term{topos}{拓扑斯}\Term{equivalence}{等价}时（于是 (5.2.2) 中所有\Term{functor}{函子}都是\Term{equivalence}{等价}）才如此。\footnote{关于 $j_{X!}$ 与\Term{topos}{拓扑斯}变换的一个交换性质，见第 XVII 讲 5.1.2。}

同样，\Term{functor}{函子} $j_{X*}$ 一般不\Term{right-exact}{右正合}，也不一定把\Term{epimorphism}{满态射}变为\Term{epimorphism}{满态射}。由这些观察可知，对于\Term{topos}{拓扑斯} $E$ 的一个\Term{object}{对象} $X$，连接 $E_{/X}$ 与 $E$ 的\Term{topos-morphism}{拓扑斯态射}的方向是毫无歧义地确定的。

\textbf{5.4。} \Term{functor}{函子} $j_X^*$ 常称为\Term{localization-functor}{局部化函子}或\Term{restriction-functor}{限制函子}。后一术语可作如下解释：把 $E$（相应地 $E_{/X}$）的\Term{object}{对象}与 $E$（相应地 $E_{/X}$）上的\Term{sheaf}{层}等同（1.2 iii)）；在此等同下，$j_{X!}$（\Term{forgetful-functor}{遗忘函子}）与 $j_X^*$ 之间的\Term{adjunction}{伴随}公式意味着，$j_X^*(F)$ 是 $E$ 上层 $F$ 在 $E_{/X}$ 上的\Term{restriction}{限制}\Term{sheaf}{层}（这里按广义理解“\Term{restriction}{限制}”，即与“包含”\Term{functor}{函子} $j_{X!}:E_{/X}\to E$ 复合）。沿用其他场合中熟悉的记号，常常也不加区别地写成
\[j_X^*(F)=F|X=F_X=\text{$F$ 在 $X$ 上的限制}.\tag{5.4.1}\]
同样，若 $E=C^\sim$，其中 $C$ 是一个（小）\Term{site}{景}，且 $X=\varepsilon(S)$，$S\in\operatorname{ob}C$，于是有 5.1 中回顾的\Term{equivalence-of-categories}{范畴等价}
\[E_{/X}=C^\sim_{/\varepsilon(S)}\xrightarrow{\approx}(C_{/S})^\sim\]
（$C_{/S}$ 配备由 $C$ 的\Term{topology}{拓扑}诱导的\Term{topology}{拓扑}），则\Term{functor}{函子} $j_X^*$ 就等同于把 $C$ 上的可变\Term{sheaf}{层} $F$ “\Term{restriction}{限制}”到\Term{category}{范畴} $C_{/S}$ 上的\Term{functor}{函子}。

这些考虑也使我们预见到，\Term{induced-topos}{诱导拓扑斯}和\Term{localization-morphism}{局部化态射} (5.2.1) 将在所有需要“在 $E$ 上\Term{local-scheme}{局部}地”论证的问题中起重要\Term{action}{作用}（见 8），实际上也就是在几乎所有涉及\Term{site}{景}或\Term{topos}{拓扑斯}的问题中起重要\Term{action}{作用}。

\textbf{5.5。} 设 $f:X\to Y$ 是 $E$ 的一个\Term{arrow}{箭头}，于是可把 $X$（更准确地说，把 $(X,f)$）视为 $E_{/Y}$ 的一个\Term{object}{对象}。显然有\Term{canonical}{典范}\Term{isomorphism}{同构}
\[(E_{/Y})_{/X}\xrightarrow{\sim}E_{/X}\tag{5.5.1}\]
\SourcePage{224}
（\Term{induced-topos}{诱导拓扑斯}的\Term{transitivity}{传递性}）。把 5.1 的构造应用于 $E_{/Y}$ 而非 $E$，便得到\Term{canonical}{典范}的\Term{topos-morphism}{拓扑斯态射}
\[\operatorname{loc}(f)\text{ 或 }f:E_{/X}\longrightarrow E_{/Y},\tag{5.5.2}\]
也称为（与 $f$ \Term{correspondence}{对应}的）\Term{localization-morphism}{局部化态射}。当 $Y$ 为 $E$ 的\Term{final-object}{终对象}时，基本上就得到 (5.2.1)。$f$ 的\Term{localization-morphism}{局部化态射}\Term{correspondence}{对应}于三个依次\Term{adjunction}{伴随}的\Term{functor}{函子}
\[f_!,\quad f^*,\quad f_*,\tag{5.5.3}\]
分别\Term{solvable}{可解}释为\Term{forgetful-functor}{遗忘函子}、\Term{restriction-functor}{限制函子}，以及一个可任选下列记号表示的\Term{functor}{函子}（$X'$ 表示自变量）：
\[\prod_{X/Y}(X'/X),\qquad\mathscr{H}om_{X/Y}(X,X'),\qquad\operatorname{Res}_{X/Y}(X'/X).\]

\textbf{5.6。} \Term{forgetful-functor}{遗忘函子}的\Term{transitivity}{传递性}说明，对 $E$ 的两个可复合\Term{morphism}{态射} $X\xrightarrow fY\xrightarrow gZ$，有\Term{topos-morphism}{拓扑斯态射}的\Term{canonical}{典范}\Term{isomorphism}{同构}
\[\operatorname{loc}(gf)\simeq\operatorname{loc}(g)\operatorname{loc}(f):E_{/X}\longrightarrow E_{/Y}\longrightarrow E_{/Z}.\]
此外，对三个可复合\Term{morphism}{态射}，上述传递\Term{isomorphism}{同构}满足通常的相容关系。因此，按照通常的用语简化，可以把 $X\mapsto E_{/X}$ 看成从 $E$ 到\Term{category}{范畴} $(V\text{-}U\text{-}\mathrm{top})$（3.3.1）的\Term{covariant-functor}{协变函子}。更准确地说，得到一个从 $E$ 到 2-\Term{category}{范畴} $(V\text{-}U\text{-}\mathrm{top})$ 的 2-\Term{functor}{函子}（通常不是严格的）。

在继续讨论\Term{induced-topos}{诱导拓扑斯}的一般性质（5.10）之前，先给几个有启发性的例子。

\textbf{5.7。} 设 $X$ 为\Term{topological-space}{拓扑空间}，因而有\Term{topos}{拓扑斯} $\operatorname{Top}(X)$（2.1）。设 $X'$ 为 $\operatorname{Top}(X)$ 的一个\Term{object}{对象}；把它解释为 $X$ 上的\Term{etale-space}{平展空间} $p:X'\to X$ 很方便。\Term{category}{范畴} $\operatorname{Top}(X)_{/X'}$ 于是等同于配备 $X$-\Term{morphism}{态射} $X''\to X'$ 的 $X$ 上\Term{etale-space}{平展空间} $X''$ 的\Term{category}{范畴}。我们知道，这样的\Term{morphism}{态射}使 $X''$ 成为 $X'$ 上的\Term{etale-space}{平展空间}，由此得到\Term{canonical}{典范}的\Term{topos}{拓扑斯}\Term{equivalence}{等价}
\[\operatorname{Top}(X)_{/X'}\xrightarrow{\approx}\operatorname{Top}(X').\]
\Term{localization-morphism}{局部化态射}于是等同于一个\Term{topos-morphism}{拓扑斯态射} $\operatorname{Top}(X')\to\operatorname{Top}(X)$，立即可见它就是结构\Term{continuous-map}{连续映射} $p:X'\to X$ 所\Term{correspondence}{对应}的\Term{morphism}{态射}
\[\operatorname{Top}(p):\operatorname{Top}(X')\longrightarrow\operatorname{Top}(X).\]
直观地说，\Term{localization-morphism}{局部化态射}不过是展映射 $p:X'\to X$ 的\Term{topos}{拓扑斯}语言表述。这一方面说明
\SourcePage{225}
“\Term{localization-morphism}{局部化态射}”这一术语是恰当的，另一方面也提醒我们谨慎使用“包含\Term{morphism}{态射}”一词；后一名称似乎主要适用于 $p$ 为\Term{open-immersion}{开浸入}的情形。因此，在 5.1 中，稳妥的做法是仅当\Term{morphism}{态射} $X\to e_E$ 为\Term{monomorphism}{单态射}，即 $X$ 可等同于 $E$ 的\Term{final-object}{终对象}的一个\Term{subobject}{子对象}时，才使用“包含\Term{morphism}{态射}”一词。

\textbf{5.8。} 设 $E$ 为\Term{topos}{拓扑斯}，$G$ 为 $E$ 中的\Term{group}{群}，$H$ 为 $G$ 的\Term{subgroup}{子群}，$X=G/H$ 为商\Term{homogeneous-space}{齐性空间}；把它视为带有左算\Term{subgroup}{子群} $G$ 的 $E$ 的\Term{object}{对象}，即\Term{classifying-topos}{分类拓扑斯} $B_G$ 的一个\Term{object}{对象}（2.4）。我们要通过定义\Term{topos}{拓扑斯}\Term{equivalence}{等价}
\[c:B_H\xrightarrow{\approx}(B_G)_{/X},\tag{5.8.1}\]
来确定\Term{induced-topos}{诱导拓扑斯} $(B_G)_{/X}$（$X=G/H$），并使下图在\Term{canonical}{典范}\Term{isomorphism}{同构}意义下交换：
\[\begin{tikzcd}B_H\arrow[r,"c","\approx"']\arrow[rr,bend right=20,"B_i"']&(B_G)_{/X}\arrow[r,"j_X"]&B_G.\end{tikzcd}\tag{5.8.2}\]
这里 $j_X$ 是 $B_G$ 中的\Term{localization-morphism}{局部化态射}，而 $B_i:B_H\to B_G$ 是由包含 $i:H\to G$ \Term{derivation}{导出}的\Term{morphism}{态射}（4.5）。

为定义 (5.8.1)，只需给出\Term{inverse-image-functor}{逆像函子} $c^*$ 和\Term{direct-image-functor}{正像函子} $c_*$ 的描述，至于验证它们确实互为\Term{quasi-inverse-functor}{拟逆}\Term{equivalence}{等价}，并给出\Term{topos-morphism}{拓扑斯态射}的\Term{commutative-diagram}{交换图} (5.8.2)，则留给读者。设 $e$ 为 $X$ 的底层 $E$-\Term{object}{对象}的\Term{subobject}{子对象}，它是由 $G$ 的\Term{unit-section}{单位截面}\Term{derivation}{导出}的 $X$ 的\Term{section}{截面}（在 $E$ 的一个选定\Term{final-object}{终对象} $e_E$ 上）的像。若 $X'$ 是 $B_G$ 中位于 $G$ 之上的一个\Term{object}{对象}，则以 $c^*(X')$ 表示 $e$ 在 $X'$ 中的\Term{inverse-image}{逆像}，
\[c^*(X')=X'\times_X e.\]
它在 $H$ 对 $X'$ 的左\Term{action}{作用}下\Term{stable}{稳定}，这一\Term{action}{作用}由 $G$ 在 $X'$ 上的给定\Term{action}{作用}\Term{restriction}{限制}而得。由此确实定义了\Term{functor}{函子}
\[c^*:(B_G)_{/X}\longrightarrow B_H.\]
另一方面，若 $X'$ 为 $B_H$ 的\Term{object}{对象}，则从 $X'$ 到 $B_H$ 的\Term{final-object}{终对象} $e_H$（即带有 $H$ 的\Term{trivial-action}{平凡作用}的 $E$ 的\Term{final-object}{终对象} $e_E$）的\Term{morphism}{态射}，经过\Term{functor}{函子} $i_!$（4.5），给出 $B_G$ 中的\Term{morphism}{态射}
\[i_!(X')\longrightarrow i_!(e_H)\simeq X,\]
\SourcePage{226}
使我们可将 $i_!(X')$ 视为 $(B_G)_{/X'}$ 的一个\Term{object}{对象}，从而得到所求\Term{functor}{函子}
\[c_*:X'\longmapsto(i_!(X')\to X):B_H\longrightarrow B_G.\]

\textbf{5.8.3。} 由 (5.8.2) 可见，若 $G$ 为某\Term{topos}{拓扑斯}中的\Term{group}{群}，$H$ 为其\Term{subgroup}{子群}，则“把\Term{operator}{算子}从 $G$ \Term{restriction}{限制}到 $H$”的\Term{functor}{函子}可以解释为\Term{localization-functor}{局部化函子}。特别地，取 $H$ 为单位\Term{subgroup}{子群}，可见“遗忘 $G$ 的\Term{action}{作用}”的\Term{functor}{函子}可以解释为\Term{localization-functor}{局部化函子}。更一般地，若 $(X,G)$ 是带有 $G$ \Term{action}{作用}的 $E$ 的一个\Term{object}{对象}（即\Term{classifying-topos}{分类拓扑斯} $B_G$ 的\Term{object}{对象}），则“遗忘 $G$ 的\Term{action}{作用}”的\Term{functor}{函子}
\[E'=(B_G)_{/(X,G)}\longrightarrow E_{/X}\]
可以解释为相对于 $E'$ 的某个适当\Term{object}{对象} $E_{G,(X,G)}$ 的\Term{localization-functor}{局部化函子}，该\Term{object}{对象}\Term{covering}{覆盖}\Term{final-object}{终对象}。只需取 $E_{G,(X,G)}=E_G\times(X,G)$，便有\Term{cartesian-diagram}{笛卡儿图}（这里 $E_G=G_S=G$，$G$ 以\Term{left-translation}{左平移}\Term{action}{作用}）
\[\begin{tikzcd}E_G\arrow[d]&E_{G,(X,G)}=Z\arrow[l]\arrow[d]\\e_G&(X,G)\arrow[l]\end{tikzcd}\tag{5.8.3.1}\]
由\Term{induced-topos}{诱导拓扑斯}的\Term{transitivity}{传递性}，\Term{induced-topos}{诱导拓扑斯} $E'=((B_G)_{X/G})_{/Z}$ \Term{equivalence}{等价}于 $((B_G)_{/E_G})_{/Z}$；而 $(B_G)_{/E_G}$ 通过 (5.8.1) 中的\Term{functor}{函子} $c^*$（取 $H=e$）\Term{equivalence}{等价}于 $E$，所以只需验证这个\Term{equivalence}{等价}把 $Z$ 变为 $E$ 的\Term{object}{对象} $X$（这是显然的），就得到所求的\Term{equivalence-of-categories}{范畴等价}
\[E'_{/Z}\xrightarrow{\approx}E_{/X}.\]
留给读者验证，此\Term{equivalence}{等价}与\Term{localization-functor}{局部化函子}
\[E'=(B_G)_{/(X,G)}\longrightarrow E'_{/Z}\]
的复合就是“遗忘 $G$ 的\Term{action}{作用}”的\Term{functor}{函子}。因此，上述\Term{cartesian-diagram}{笛卡儿图}通过取\Term{induced-topos}{诱导拓扑斯}，给出\Term{topos-morphism}{拓扑斯态射}的图
\[\begin{tikzcd}(B_G)_{/E_G}\approx E\arrow[d,"j"']&(B_G)_{/Z}\approx E'_{/Z}\approx E_{/X}\arrow[l,"i"']\arrow[d,"j'"]\\B_G&(B_G)_{/(X,G)}=E'\arrow[l,"f"]\end{tikzcd}\tag{5.8.3.2}\]
它在\Term{canonical}{典范}\Term{isomorphism}{同构}意义下交换（5.6）。其中，竖直\Term{arrow}{箭头} $j,j'$ \Term{correspondence}{对应}的\Term{inverse-image-functor}{逆像函子}均为“遗忘 $G$ 的\Term{action}{作用}”的\Term{functor}{函子}，而\Term{inverse-image-functor}{逆像函子} $i^*$ 等同于\Term{localization-functor}{局部化函子} $E\to E_{/X}$。
\SourcePage{227}
这个图还是下述 5.11 意义下的“2-\Term{cartesian-diagram}{笛卡儿图}”。

\textbf{5.8.4。} 留给读者把上述考虑推广到\Term{pro-group}{pro-群} $G=(G_i)$ 的情形，其中只考虑 $H$ 定义为 $G\to G_{i_0}$ 的“\Term{kernel}{核}”的特殊情形。（对\Term{profinite-group}{投有限群}而言，这就是说，只考虑 $G$ 的开\Term{subgroup}{子群} $H$，即闭且指数\Term{finite}{有限}的\Term{subgroup}{子群}。）

\begin{SGAresult}{习题}{5.9}
设 $E$ 为\Term{topos}{拓扑斯}，$G$ 为 $E$ 中的\Term{group}{群}，$B_G$ 为 $G$ 的\Term{classifying-topos}{分类拓扑斯}（2.4），$\pi:B_G\to E$ 为由 $G$ 到 $E$ 的\Term{unit-group}{单位群} $e$ 的\Term{homomorphism}{同态}\Term{derivation}{导出}的\Term{topos-morphism}{拓扑斯态射}（这里利用了 $B_e\approx E$）。

 a) 设 $E_G=G_s$ 为 $B_G$ 的\Term{object}{对象}，其底层 $E$-\Term{object}{对象}是 $G$，$G$ 的\Term{action}{作用}由\Term{left-translation}{左平移}定义。把 $\pi^*(G)$ 看成 $B_G$ 中的\Term{group}{群}（即 $E$ 中的\Term{group}{群} $G$，其上配备 $G$ 的\Term{trivial-action}{平凡作用}）。证明，由 $G$ 在 $G_s$ 上的右\Term{translation}{平移}定义的 $E$-\Term{morphism}{态射} $G_S\times G\to G_s$，与 $G$ 在 $G_s$ 及 $G=\pi^*(G)$ 上的\Term{action}{作用}相容，并定义 $B_G$ 的\Term{morphism}{态射}
\[E_G\times\pi^*(G)\longrightarrow E_G.\]
证明，此\Term{morphism}{态射}使 $E_G$ 成为带有 $B_G$-\Term{group}{群} $\pi^*(G)$ 的\Term{right-action}{右作用}的 $B_G$-\Term{object}{对象}，且该\Term{action}{作用}使 $E_G$ 成为 $\pi^*(G)$ 下的\Term{torsor}{挠子}（即 $B_G$ 的\Term{final-object}{终对象}可以由\Term{object}{对象} $X_i$ \Term{covering}{覆盖}，使得对每个 $i$，$E_G$ 在 $X_i$ 上的\Term{restriction}{限制}\Term{isomorphism}{同构}于\Term{trivial-bundle}{平凡}\Term{operator}{算子}丛 $(\pi^*(G)_{X_i})_d$）。

 b) 设 $q:E'\to E$ 为 $E$ 上的\Term{topos}{拓扑斯}。对于 $E$ 上的任意\Term{topos}{拓扑斯} $B$，定义\Term{category}{范畴} $\mathscr{H}om\mathrm{top}_E(E',B)$：其\Term{object}{对象}为从 $E'$ 到 $B$ 的\Term{topos-morphism}{拓扑斯态射}，并在给定\Term{isomorphism}{同构}意义下与结构\Term{morphism}{态射} $E'\to E$ 和 $B\to E$ 相容。取 $B=B_G$，用公式 $f\mapsto f^*(E_G)$ 定义\Term{equivalence-of-categories}{范畴等价}
\[\mathscr{H}om\mathrm{top}_E(E',B_G)\longrightarrow\operatorname{Tors}(E,q^*(G)),\]
其中右边表示 $E$ 上（右）$q^*(G)$-\Term{torsor}{挠子}的\Term{category}{范畴}。

 c) 设 $H$ 为 $G$ 的\Term{subgroup}{子群}，于是 $\pi^*(H)$ 是 $p^*(G)$ 的\Term{subgroup}{子群}，并通过\Term{restriction}{限制}算\Term{subgroup}{子群}而\Term{action}{作用}于 $E_G$，从而有\Term{quotient-object}{商对象} $X=E_G/\pi^*(H)$，
\SourcePage{228}
它正是带有 $G$ 的通常左\Term{action}{作用}的\Term{object}{对象} $G/H$。设 $e_G$ 为 $B_G$ 的\Term{final-object}{终对象}（即带有 $G$ 的\Term{trivial-action}{平凡作用}——这里别无选择——的 $E$ 的\Term{final-object}{终对象} $e_E$），考虑 $B_G$ 中的\Term{morphism}{态射}图
\[\begin{tikzcd}E_G\arrow[dr]\arrow[dd]\\&X=E_G/\pi^*(H)\arrow[dl]\\e_G.\end{tikzcd}\]
在这些\Term{object}{对象}上取由 $B_G$ 诱导的\Term{topos}{拓扑斯}，便得到一个在\Term{canonical}{典范}\Term{isomorphism}{同构}意义下交换的\Term{topos-morphism}{拓扑斯态射}图（5.6）。证明，该图\Term{equivalence}{等价}于
\[\begin{tikzcd}E=B_e\arrow[dr]\arrow[dd]\\&B_H\arrow[dl]\\B_G,\end{tikzcd}\]
其中三个\Term{arrow}{箭头}均为依 4.5 与 $E$ 中\Term{group}{群}\Term{morphism}{态射} $e\to H\to G$ \Term{correspondence}{对应}的\Term{arrow}{箭头}。（因而，直观上可将\Term{classifying-topos}{分类拓扑斯} $B_H$ 解释为 $B_G$ 上的一个\Term{homogeneous-space}{齐性空间}，其\Term{group}{群}为 $\pi^*(G)$，\Term{correspondence}{对应}于 $B_G$ 上的\Term{universal}{普遍}\Term{torsor}{挠子}（即\Term{principal-homogeneous-bundle}{主齐性丛}）$E_G$。）

 d) 在 $X'$ 为 $X$ 的\Term{etale}{平展}\Term{covering}{覆盖}（\Term{locally-trivial}{局部平凡}，见 2.7.4）的情形，重新讨论例 5.7，并用本习题中的考虑解释它。
\end{SGAresult}

\textbf{5.10。\Term{induced-topos}{诱导拓扑斯}的\Term{inverse-image}{逆像}。} 设 $f:E'\to E$ 为 $U$-\Term{topos-morphism}{拓扑斯态射}，$X$ 为 $E$ 的\Term{object}{对象}，$X'=f^*(X)$ 为它在 $E$ 中的\Term{inverse-image}{逆像}。我们将定义一个\Term{topos-morphism}{拓扑斯态射}图
\[\begin{tikzcd}E_{/X}\arrow[r,"f_{/X}"]\arrow[d,"j_X"']&E'_{/X'}\arrow[d,"j_{X'}"]\\E&E'\arrow[l,"f"]\end{tikzcd}\tag{5.10.1}\]
其中 $j_X,j_{X'}$ 为\Term{localization-morphism}{局部化态射}（5.1），$f_{/X}$ 将在下面定义，
\SourcePage{229}
且此图在\Term{canonical}{典范}\Term{isomorphism}{同构}意义下交换。这里通过\Term{inverse-image-functor}{逆像函子}
\[(f_{/X})^*:E_{/X}\longrightarrow E'_{/X'}\tag{5.10.2}\]
定义 $f_{/X}$；取它为由 $f^*$ 在显然的意义下“诱导”的\Term{functor}{函子}。因为包含\Term{functor}{函子} $j_{X!}:E_{/X}\to E$ 和 $j_{X'!}:E_{/X}\to E$ 与 $U$-\Term{inductive-limit}{归纳极限}及\Term{fiber-product}{纤维积}交换，并且是\Term{conservative-functor}{保守}的，所以\Term{functor}{函子} $(f_{/X})^*$ 同诱导它的\Term{functor}{函子} $f^*$ 一样，与 $U$-\Term{inductive-limit}{归纳极限}及\Term{fiber-product}{纤维积}交换；此外 $(f_{/X})^*$ 显然把\Term{final-object}{终对象} $X$ 变为\Term{final-object}{终对象} $X'$，故它\Term{left-exact}{左正合}，从而\Term{correspondence}{对应}于一个\Term{topos-morphism}{拓扑斯态射}
\[f_{/X}:E'_{/X'}\longrightarrow E_{/X},\tag{5.10.3}\]
该\Term{morphism}{态射}在唯一\Term{isomorphism}{同构}意义下确定（3.1.1）。另一方面，由构造及 $f^*$ 与积 $X\times Y$ 交换这一事实，从 (5.10.1) 取\Term{inverse-image-functor}{逆像函子}所得的图在\Term{canonical}{典范}\Term{isomorphism}{同构}意义下交换，因此 (5.10.1) 本身也在一个唯一\Term{isomorphism}{同构}意义下交换，这个\Term{isomorphism}{同构}在两个复合 $f\circ j_{X'}$ 和 $j_X\circ f_{/X}$ 的\Term{inverse-image-functor}{逆像函子}上诱导该\Term{canonical}{典范}\Term{isomorphism}{同构}（3.2.1）。

\textbf{5.10.4。} 当 $f:E'\to E$ \Term{correspondence}{对应}于\Term{topological-space}{拓扑空间}的\Term{continuous-map}{连续映射} $Y'\to Y$（4.1）时，$X$ 等同于 $Y$ 上的一个\Term{etale-space}{平展空间}，而 $X'$ 等同于\Term{fiber-product}{纤维积} $X\times_Y Y'$；把\Term{induced-topos}{诱导拓扑斯} $\operatorname{Top}(Y)_{/X}$ 和 $\operatorname{Top}(Y')_{/X'}$ 分别等同于 $\operatorname{Top}(X)$ 和 $\operatorname{Top}(X')$（5.7），则刚构造的\Term{morphism}{态射} $f_{/X}$ 就是（在唯一\Term{isomorphism}{同构}意义下）由\Term{canonical}{典范}\Term{continuous-map}{连续映射} $X'\to Y'$ \Term{derivation}{导出}的\Term{morphism}{态射} $\operatorname{Top}(X')\to\operatorname{Top}(Y')$。借助这个例子，可以说\Term{topos}{拓扑斯} $E'_{/X'}$ 起着 $E_{/X}$ 与 $E'$ 在 $E$ 上的\Term{fiber-product}{纤维积}的\Term{action}{作用}。下述结果使这一直观更加精确。

\begin{SGAresult}{命题}{5.11}
沿用 5.10 的记号，图 (5.10.1) 连同相容\Term{isomorphism}{同构} $\alpha:f\circ j_{X'}\to j_X\circ f_{/X}$ 是“2-\Term{cartesian}{笛卡儿}”的。确切地说，对任意 $U$-\Term{topos}{拓扑斯} $F$，若给每个\Term{topos-morphism}{拓扑斯态射} $g:F\to E'_{/X'}$ 配以三元组
\[(f_{/X}\circ g,j_{X'}\circ g,\alpha*g)=(g_1,g_2,\beta),\]
其中 $g_1:F\to E_{/X}$ 和 $g_2:F\to E'$ 为\Term{topos-morphism}{拓扑斯态射}，而 $\beta:j_X\circ g_1\to f\circ g_2$ 是从 $F$ 到 $E$ 的\Term{topos-morphism}{拓扑斯态射}的\Term{isomorphism}{同构}，则得到 $\mathscr{H}om\mathrm{top}(F,E'_{/X})$ 与所有上述三元组 $(g_1,g_2,\beta)$ 构成的\Term{category}{范畴}
\[\mathscr{H}om\mathrm{top}(F,E_{/X})\mathbin{\mathop{\times}\limits^{2}_{\mathscr{H}om\mathrm{top}(F,E)}}\mathscr{H}om\mathrm{top}(F,E')\]
之间的\Term{equivalence-of-categories}{范畴等价}。
\end{SGAresult}
（注意：在\Term{cartesian-product}{笛卡儿积}\Term{signature}{符号}上方加上 $2$，是为提醒读者这不是\Term{category}{范畴}的普通\Term{fiber-product}{纤维积}，而是命题中明确说明的意义下的“2-\Term{fiber-product}{纤维积}”。）

\SourcePage{230}
为证明 5.11，我们将借助\Term{correspondence}{对应}的\Term{inverse-image-functor}{逆像函子}来具体表述\Term{topos-morphism}{拓扑斯态射}。为此需要下面 5.11.1 至 5.12 中给出的辅助结果。

\textbf{5.11.1。} 一般地，考虑如下情形：给定两个\Term{category}{范畴} $E,E'$，$E$ 的一个可作积的\Term{object}{对象} $X$，即\Term{functor}{函子}
\[j:A\longmapsto A\times X:E\longrightarrow E_{/X}\]
有定义，以及一个\Term{left-exact}{左正合}\Term{functor}{函子}
\[\varphi:E_{/X}\longrightarrow E'.\tag{*}\]
因为 $E_{/X}$ 有\Term{final-object}{终对象} $(X,\operatorname{id}_X)=X$，故 $E'$ 也如此，且其\Term{final-object}{终对象}为 $e'\simeq\varphi(X)$。

选定 $E'$ 的\Term{final-object}{终对象} $e'$ 后，我们要对上述每个 $\varphi$ 配以二元组
\[(\Psi,u),\quad\Psi=\varphi\circ j:E\to E',\quad u\in\operatorname{Hom}(e',\Psi(X)).\tag{**}\]
为定义 $u$，注意 $\Psi(X)=X\times X$ 在 $E_{/X}$ 的\Term{final-object}{终对象} $X$ 上有一个\Term{canonical}{典范}\Term{section}{截面} $\delta_X$，即\Term{diagonal-section}{对角截面}，并定义
\[u=\varphi(\delta_X):\varphi(X)=e'\longrightarrow\varphi(X\times X)=\Psi(X).\]
我们断言，由二元组 (**) 可在唯一\Term{isomorphism}{同构}意义下重构 (*) 中的\Term{functor}{函子} $\varphi$。为此，对 $E_{/X}$ 的任意\Term{object}{对象} $p:X'\to X$，考虑 $E_{/X}$ 中的\Term{cartesian-diagram}{笛卡儿图}
\[\begin{tikzcd}X'\arrow[r]\arrow[d]&j(X')=X'\times X\arrow[d,"j(p)"]\\X\arrow[r,"\delta_X"]&j(X)=X\times X,\end{tikzcd}\]
其中上方水平\Term{arrow}{箭头}为\Term{graph-morphism}{图态射} $\Gamma_p=(\operatorname{id}_X,p)$。对此图应用\Term{left-exact}{左正合}\Term{functor}{函子} $\varphi$，并利用 $\Psi=\varphi\circ j$ 的定义，
\SourcePage{231}
得到\Term{cartesian-diagram}{笛卡儿图}
\[\begin{tikzcd}\varphi(X',p)\arrow[r]\arrow[d]&\Psi(X')\arrow[d,"\Psi(p)"]\\e'\arrow[r]&\Psi(X).\end{tikzcd}\]
换言之，得到\Term{canonical}{典范}\Term{isomorphism}{同构}
\[\varphi(X',p)\simeq\Psi(X')\times_{\Psi(X)}e'.\tag{5.11.1.1}\]
它显然关于 $E_{/X}$ 的\Term{object}{对象} $(X',p)$ 是\Term{functor}{函子}性的，这就具体说明了如何借二元组 (**) 中的 $(\Psi,u)$，在\Term{isomorphism}{同构}意义下重构\Term{functor}{函子} $\varphi:E_{/X}\to E'$。另外，\Term{functor}{函子} $\Psi$ 是\Term{left-exact}{左正合}的；更一般地，因为 $j$ 与\Term{projective-limit}{投射极限}交换，故 $\Psi$ 与 $\varphi$ 所交换的任何类型的\Term{projective-limit}{投射极限}都交换。

反过来，现在假定 $E,E'$ 中\Term{finite}{有限}\Term{projective-limit}{投射极限}均\Term{representable}{可表}，并从一个二元组
\[(\Psi,u),\quad\Psi:E\longrightarrow E',\quad u\in\operatorname{Hom}(e',\Psi(X))\]
出发，其中\Term{functor}{函子} $\Psi$ \Term{left-exact}{左正合}，便可通过 (5.11.1.1) 定义\Term{functor}{函子} $\varphi:E_{/X}\to E'$。立即可验证，此\Term{functor}{函子}\Term{left-exact}{左正合}；更一般地，它与 $E$ 中\Term{representable}{可表}且 $\Psi$ 与之交换的任何类型的\Term{projective-limit}{投射极限}都交换。这立即由连接 $E$ 和 $E_{/X}$ 中所计算的\Term{projective-limit}{投射极限}的\Term{cartesian-diagram}{笛卡儿图}
\[\begin{tikzcd}\varprojlim_I^{E}X'_i\arrow[d]&\varprojlim_I^{E_{/X}}X'_i\arrow[l]\arrow[d]\\\varprojlim_I^{E}X&X\arrow[l,"\delta"']\end{tikzcd}\]
得出：其中 $\delta$ 表示到常值为 $X$ 的\Term{functor}{函子}的\Term{projective-limit}{投射极限}的\Term{diagonal-morphism}{对角态射}。对此图应用\Term{functor}{函子} $\Psi$，得到 $E'$ 中的\Term{cartesian-diagram}{笛卡儿图}，再与从 $\Psi(X)\leftarrow e'$ \Term{derivation}{导出}的\Term{cartesian-square}{笛卡儿方块}复合，便得到一个复合\Term{cartesian-square}{笛卡儿方块}，表达出 $\varphi$ 与所考虑的\Term{projective-limit}{投射极限}的相容性。

可由 $\varphi$ 在\Term{isomorphism}{同构}意义下重构二元组 $(\Psi,u)$。事实上，有\Term{canonical}{典范}\Term{isomorphism}{同构}
\[\Psi(X')=\varphi(j(X')),\tag{5.11.1.2}\]
它由 (5.11.1.1) 中把 $(X',p)$ 换成 $j(X')=(X'\times X,\operatorname{pr}_2)$，并注意到 $\Psi(X'\times X)\simeq\Psi(X')\times\Psi(X)$ 而得。上述\Term{isomorphism}{同构}显然关于 $X'$ \Term{functor}{函子}性，
\SourcePage{232}
即它给出\Term{isomorphism}{同构} $\Psi\simeq\varphi\circ j$；同样可验证，在此\Term{isomorphism}{同构}下 $u:e'\to\Psi(X)$ 等同于 $\varphi(\delta_X)$。由此已得到下述引理的主要内容。

\begin{SGAresult}{引理}{5.11.2}
设 $E,E'$ 为\Term{finite}{有限}\Term{projective-limit}{投射极限}均\Term{representable}{可表}的两个\Term{category}{范畴}，$X$ 为 $E$ 的\Term{object}{对象}，$\mathscr{S}\!\operatorname{ex}(E,E')$ 为从 $E$ 到 $E'$ 的\Term{left-exact}{左正合}\Term{functor}{函子} $\Psi$ 的\Term{category}{范畴}，$\mathscr{S}\!\operatorname{ex}(E_{/X},E')$ 为类似地定义的从 $E_{/X}$ 到 $E'$ 的\Term{left-exact}{左正合}\Term{functor}{函子} $\varphi$ 的\Term{category}{范畴}。则\Term{category}{范畴} $\mathscr{S}\!\operatorname{ex}(E_{/X},E')$ 与二元组 $(\Psi,u)$ 的\Term{category}{范畴} $\mathscr{S}\!\operatorname{ex}(E,E')_{/\Gamma}$ \Term{equivalence}{等价}，其中 $\Psi\in\operatorname{ob}\mathscr{S}\!\operatorname{ex}(E,E')$，$u:e'\to\Psi(X)$，而 $e'$ 是 $E'$ 的\Term{final-object}{终对象}。该\Term{equivalence}{等价}及其\Term{quasi-inverse-functor}{拟逆}\Term{functor}{函子}的定义已在 5.11.1 中明确给出。若 $\varphi$ 与 $(\Psi,u)$ 相\Term{correspondence}{对应}，则 $\varphi$ 与某个指定类型的\Term{projective-limit}{投射极限}交换，当且仅当 $\Psi$ 也如此。若 $E,E'$ 中\Term{base-change-functor}{换基函子}与某个指定类型的\Term{inductive-limit}{归纳极限}交换，则 $\varphi$ 与该类型的\Term{inductive-limit}{归纳极限}交换，当且仅当 $\Psi$ 也如此。
\end{SGAresult}

留给读者具体说明二元组 $(\Psi,u)$ \Term{correspondence}{对应}的\Term{category}{范畴} $\mathscr{S}\!\operatorname{ex}(E,E')_{/\Gamma}$ 的结构。定义\Term{functor}{函子}
\[\Gamma:\mathscr{S}\!\operatorname{ex}(E,E')\to(\mathrm{Ens}),\qquad\Gamma(\Psi)=\operatorname{Hom}(e',\Psi(X)),\]
则二元组 $(\psi,u)$ 中的 $u$ \Term{solvable}{可解}释为 $\Gamma(\Psi)$ 的一个元素，即 $\mathscr{S}\!\operatorname{ex}(E,E')^\circ$ 上的\Term{presheaf}{预层}\Term{category}{范畴}中的一个\Term{homomorphism}{同态} $\Psi\to\Gamma$；这说明了 5.11.1 的陈述中所用记号 $/\Gamma$ 的理由。为证明 5.11.2，还需在 5.11.1 中具体说明所考虑构造关于变化的 $\varphi$ 或 $(\psi,u)$ 的\Term{functor}{函子}性，这基本上是显然的，留给读者；最后还须验证与\Term{inductive-limit}{归纳极限}交换的断言，借助连接这些\Term{functor}{函子}的显式公式 (5.11.1.1) 和 (5.11.1.2)，这一点也显然。

特别地，当 $E,E'$ 为 $U$-\Term{topos}{拓扑斯}，并借助\Term{correspondence}{对应}的\Term{inverse-image-functor}{逆像函子}来解释\Term{topos-morphism}{拓扑斯态射}时，可得：
\begin{SGAresult}{命题}{5.12}
设 $E,E'$ 为两个 $U$-\Term{topos}{拓扑斯}，$X$ 为 $E$ 的一个\Term{object}{对象}，并考虑 $\mathscr{H}om\mathrm{top}(E',E)$ 上的\Term{contravariant-functor}{反变函子}
\[\gamma:f\longmapsto\Gamma(E',f^*(X))=\operatorname{Hom}(e',f^*(X)):\mathscr{H}om\mathrm{top}(E',E)^\circ\longrightarrow(\mathrm{Ens}).\]
则有\Term{equivalence-of-categories}{范畴等价}
\[\mathscr{H}om\mathrm{top}(E',E_{/X})\xrightarrow{\approx}\mathscr{H}om\mathrm{top}(E',E)_{/\gamma},\tag{5.12.1}\]
其中右边是 $\mathscr{H}om\mathrm{top}(E',E)^{\wedge}_{/\gamma}$ 中由二元组 $(f,u)$ 构成的\Term{full-subcategory}{满子范畴}，这里 $f\in\operatorname{Homtop}(E',E)$ 且 $u\in\gamma(f)$，即 $u\in\Gamma(E',f^*(X))$。此\Term{functor}{函子}把任意\Term{topos-morphism}{拓扑斯态射} $h:E'\to E_{/X}$ 映为二元组 $(f,u)$，其中 $f$ 是复合
\[f=j_X\circ h:E'\longrightarrow E_{/X}\longrightarrow E,\]
\SourcePage{233}
而\Term{morphism}{态射}
\[u:e'=h^*(X,\operatorname{id}_X)\longrightarrow f^*(X)=h^*(j_X^*(X))=h^*(X\times X)\]
是对\Term{diagonal-morphism}{对角态射} $\delta_X:X\to X\times X$ 应用 $h^*$ 而得。
\end{SGAresult}

\textbf{5.12.2。} 利用 5.12，5.11 的证明就几乎显然了，实质上归结为下述考虑（它可以替代一个并不能带来更多启发的“形式化”证明）：给定一个\Term{topos-morphism}{拓扑斯态射} $g:F\to E'_{/X'}$，“等同于”给定\Term{topos-morphism}{拓扑斯态射} $g_2:F\to E'$ 及 $g_2^*(X')$ 的一个\Term{section}{截面} $u$；而 $g_2^*(X')=g_2^*(f^*(X))=(fg_2)^*(X)$，因此给定 $u$ 也\Term{equivalence}{等价}于给定\Term{topos-morphism}{拓扑斯态射} $fg_2:F\to E$ 到\Term{topos-morphism}{拓扑斯态射} $g_1:F\to E_{/X}$ 的一个\Term{lifting}{提升}。这正说明了给定\Term{topos-morphism}{拓扑斯态射} $g:F\to E'_{/X'}$，实质上\Term{equivalence}{等价}于给定 5.11 中那样的三元组 $(g_1,g_2,\alpha)$。

\textbf{注记 5.13。} 我们将看到（第 15 节），对任意\Term{topos}{拓扑斯}图
\[\begin{tikzcd}E_1\arrow[d]\\E&E_2\arrow[l]\end{tikzcd}\]
都存在属于 $V$ 的 $U$-\Term{topos}{拓扑斯}的 2-\Term{category}{范畴} $(V\text{-}U\text{-}\mathrm{Top})$ 意义下的 2-\Term{fiber-product}{纤维积}（在\Term{equivalence}{等价}意义下与满足 $U\in V$ 的\Term{universe}{宇宙} $V$ 的选择无关）。关于除 5.11 以外“\Term{topos}{拓扑斯}的\Term{fiber-product}{纤维积}”的其他自然例子，见 Giraud 的书 [3] 的最后一章。

\begin{SGAresult}{习题}{5.14}
 a) 设 $F,G$ 是\Term{topos}{拓扑斯} $E$ 上的两个\Term{topos}{拓扑斯}。定义二元组 $(f,\alpha)$ 的\Term{category}{范畴} $\mathscr{H}om\mathrm{top}_E(F,G)$，其中 $f:F\to G$ 是\Term{topos-morphism}{拓扑斯态射}，而 $\alpha:p\to q\circ f$ 是从 $E$ 到 $G$ 的\Term{topos-morphism}{拓扑斯态射}的\Term{isomorphism}{同构}（$p:F\to E$ 和 $q:G\to E$ 为结构\Term{topos-morphism}{拓扑斯态射}）。定义\Term{pairing}{配对}\Term{functor}{函子}
\[\mathscr{H}om\mathrm{top}_E(F,G)\times\mathscr{H}om\mathrm{top}_E(G,H)\longrightarrow\mathscr{H}om\mathrm{top}_E(F,H)\]
及结合律\Term{isomorphism}{同构}。

 b) 设 $X,Y$ 为\Term{topos}{拓扑斯} $E$ 的两个\Term{object}{对象}。定义一个从\Term{set}{集合} $\operatorname{Hom}(X,Y)$ 所定义的\Term{discrete-category}{离散范畴}到\Term{category}{范畴} $\mathscr{H}om\mathrm{top}_E(E_{/X},E_{/Y})$ 的自然\Term{functor}{函子}。说明它与 $E$ 中\Term{morphism}{态射}的复合及 a) 中所考虑\Term{pairing}{配对}的相容性。

 c) 证明 b) 中的\Term{functor}{函子}为\Term{equivalence-of-categories}{范畴等价}。特别地，若把\Term{induced-topos}{诱导拓扑斯} $E_{/X}$ 作为 $E$ 上的\Term{topos}{拓扑斯}，已知它“到 $E$-\Term{equivalence}{等价}为止”的信息，
\SourcePage{234}
则\Term{topos}{拓扑斯} $E$ 的\Term{object}{对象} $X$ 可以在唯一\Term{isomorphism}{同构}意义下重构。\footnote{此结果归于 P. Deligne。$E$ 为\Term{punctual-topos}{点拓扑斯}的情形先前已由 Hakim 女士处理。}
\end{SGAresult}

\SGAsection{6}{拓扑斯的点与纤维函子}
\textbf{6.0。} 设 $P$ 为标准的 $U$-\Term{punctual-topos}{点拓扑斯}（2.2）
\[P=(U\text{-}\mathrm{Ens}).\]
\Term{signature}{符号} $P$ 表示一个性质如同点的几何\Term{object}{对象}，而 $(U\text{-}\mathrm{Ens})$ 表示一个“大\Term{category}{范畴}”，解释为“$P$ 上的\Term{sheaf}{层}\Term{category}{范畴}”；两者所附带的直观颇有不同。因此，虽然从严格逻辑上说它们表示同一个\Term{object}{对象}，仍宜依上下文只使用其中一个记号。

\begin{SGAresult}{定义}{6.1}
设 $E$ 为 $U$-\Term{topos}{拓扑斯}。从\Term{punctual-topos}{点拓扑斯} $P$（6.0）到 $E$ 的任意\Term{topos-morphism}{拓扑斯态射} $P\to E$ 称为 $E$ 的\Term{point-topos}{点}。\Term{category}{范畴} $\mathscr{H}om\mathrm{top}(P,E)$（3.2）称为 $E$ 的点\Term{category}{范畴}，记作 $\operatorname{Points}(E)$ 或 $\operatorname{Pt}(E)$。若 $p:P\to E$ 是 $E$ 的一个点，\Term{correspondence}{对应}于\Term{inverse-image-functor}{逆像函子} $p^*:E\to(U\text{-}\mathrm{Ens})$，而 $F$ 是 $E$ 的\Term{object}{对象}，则\Term{set}{集合} $p^*(F)$ 称为 $E$ 在 $p$ 处的\Term{fiber}{纤维}，记为 $F_p$。
\end{SGAresult}

\textbf{6.1.1。} 当然，若 $F$ 为 $E$ 的\Term{group-object}{群对象}（或\Term{ring}{环}\Term{object}{对象}，或……），则它在 $p$ 处的\Term{fiber}{纤维} $F_p$ 是一个\Term{group}{群}（或\Term{ring}{环}，或……）（3.1.2）。

\textbf{6.1.2。} 由 3.2.1，$E$ 的点\Term{category}{范畴}通过\Term{functor}{函子} $p\mapsto p^*$ \Term{equivalence}{等价}于如下\Term{functor}{函子}的\Term{category}{范畴}的\Term{opposite-category}{反范畴}：
\[\varphi:E\longrightarrow(U\text{-}\mathrm{Ens}),\]
它们与 $U$-\Term{inductive-limit}{归纳极限}交换且\Term{left-exact}{左正合}。因此，给定 $E$ 的一个点，实质上\Term{equivalence}{等价}于给定这样的\Term{functor}{函子} $\varphi$。

\begin{SGAresult}{定义}{6.2}
与 $U$-\Term{inductive-limit}{归纳极限}交换且\Term{left-exact}{左正合}的任意\Term{functor}{函子} $\varphi:E\to(U\text{-}\mathrm{Ens})$ 称为 $U$-\Term{topos}{拓扑斯} $E$ 的\Term{fiber-functor}{纤维函子}。由 $E$ 上的\Term{fiber-functor}{纤维函子}构成的 $\mathscr{H}om(E,(U\text{-}\mathrm{Ens}))$ 的\Term{full-subcategory}{满子范畴}，称为 $E$ 上的\Term{fiber-functor}{纤维函子}\Term{category}{范畴}，记为 $\operatorname{Fib}(E)$。
\end{SGAresult}

\textbf{6.2.3。} 由 6.1.2，$E$ 上的\Term{fiber-functor}{纤维函子}\Term{category}{范畴} $\operatorname{Fib}(E)$ 通过从 $E$ 的点\Term{category}{范畴} $\operatorname{Point}(E)$ 到它的\Term{functor}{函子} $p\mapsto p^*$，\Term{equivalence}{等价}于后者的\Term{opposite-category}{反范畴}：
\[\operatorname{Point}(E)\xrightarrow{\approx}\operatorname{Fib}(E)^\circ.\tag{6.2.1.1}\]
\SourcePage{235}
\Term{functor}{函子} $\varphi:E\to(U\text{-}\mathrm{Ens})$ 是\Term{fiber-functor}{纤维函子}的充要条件是，它为 $E$ 的某个点 $p$ 所\Term{correspondence}{对应}的\Term{fiber-functor}{纤维函子} $F\mapsto F_p$（6.1）。换言之，\Term{functor}{函子} (6.2.1.1) 是\Term{surjection}{满射}的。另外，\Term{functor}{函子} $\varphi$ 是\Term{fiber-functor}{纤维函子}的充要条件是，它\Term{left-exact}{左正合}且把\Term{covering-family}{覆盖族}变为\Term{surjective-family}{满射族}（1.7）。后一条件也\Term{representable}{可表}述为：$\varphi$ 与 $U$-和交换，并把\Term{epimorphism}{满态射}变为\Term{epimorphism}{满态射}。

\textbf{6.3。} 设 $C\in U$ 为\Term{site}{景}。$C^\sim$ 的任意点称为\Term{site}{景} $C$ 的点；\Term{site}{景} $C$ 的点\Term{category}{范畴}定义为
\[\operatorname{Point}(C)=\operatorname{Point}(C^\sim).\]
另一方面，考虑\Term{functor}{函子}
\[\varphi\longmapsto\varphi|C=\varphi\circ\varepsilon_C:\operatorname{Fib}(C^\sim)\longrightarrow\mathscr{H}om(C,(U\text{-}\mathrm{Ens})).\tag{6.3.1}\]
它是\Term{fully-faithful-functor}{全忠实}的，其\Term{essential-image}{本质像}为\Term{full-subcategory}{满子范畴} $\operatorname{Morsite}((U\text{-}\mathrm{Ens}),C)^\circ$（4.9.4），也将它记为 $\operatorname{Fib}(C)$。\Term{functor}{函子} $\Psi:C\to(U\text{-}\mathrm{Ens})$ 若属于刚明确说明的 (6.3.1) 的\Term{essential-image}{本质像}，就称为\Term{site}{景} $C$ 上的\Term{fiber-functor}{纤维函子}。这样的\Term{functor}{函子}因而是\Term{continuous}{连续}的（从而由第 III 讲 1.6，它把\Term{covering-family}{覆盖族}变为\Term{covering-family}{覆盖族}，即\Term{surjective-family}{满射族}），并且\Term{left-exact}{左正合}。当 $C$ 中的\Term{finite}{有限}\Term{projective-limit}{投射极限}均\Term{representable}{可表}时（实际遇到的情况几乎全部满足这个条件），上述性质刻画了 $C$ 上的\Term{fiber-functor}{纤维函子}：它们恰是那些\Term{left-exact}{左正合}且把\Term{covering-family}{覆盖族}变为\Term{surjective-family}{满射族}的\Term{functor}{函子} $\Psi:C\to(U\text{-}\mathrm{Ens})$（4.9.4）。

应当记住，\Term{functor}{函子} $\varphi\mapsto\varphi|C$ 是\Term{equivalence-of-categories}{范畴等价}
\[\operatorname{Fib}(C^\sim)\xrightarrow{\approx}\operatorname{Fib}(C),\]
因此给定\Term{topos}{拓扑斯} $C^\sim$ 上的\Term{fiber-functor}{纤维函子}（或者这个\Term{topos}{拓扑斯}的一个点），实质上\Term{equivalence}{等价}于给定\Term{site}{景} $C$ 上的\Term{fiber-functor}{纤维函子}。给定 $C$ 上的\Term{fiber-functor}{纤维函子} $\Psi$，它因而在\Term{isomorphism}{同构}意义下来自 $C^\sim$ 上的\Term{fiber-functor}{纤维函子} $\varphi$，则可借 $\Psi$ 由公式
\[\varphi(F)\simeq\varinjlim_{C_{/F}}\Psi(X),\tag{6.3.2}\]
在\Term{canonical}{典范}\Term{isomorphism}{同构}意义下重构 $\varphi$。这里 $C_{/F}$ 为配备\Term{morphism}{态射} $X\to F$（在 $\widehat C$ 中），即配备 $F(X)$ 的一个元素的 $C$ 的\Term{object}{对象} $X$ 所构成的\Term{category}{范畴}。公式 (6.3.2) 立即由
\SourcePage{236}
$\varphi$ 与\Term{inductive-limit}{归纳极限}交换，及关于 $F$ \Term{functor}{函子}性的\Term{canonical}{典范}\Term{isomorphism}{同构}（第 II 讲 4.1.1）
\[F\simeq\varinjlim_{C_{/F}}\varepsilon_C(X)\]
得出。更一般地，当 $G$ 为 $C$ 上的\Term{presheaf}{预层}时，有关于 $G$ \Term{functor}{函子}性的\Term{canonical}{典范}\Term{isomorphism}{同构}
\[\varphi(\underline G)\simeq\varinjlim_{C_{/G}}\Psi(X),\tag{6.3.3}\]
其中 $\underline G$ 为 $G$ 的\Term{associated-sheaf}{伴随层}（第 II 讲 4.1.1）。事实上，我们知道公式
\[\underline G\simeq\varinjlim_{C_{/G}}\varepsilon_C(X).\]

\textbf{6.4.0。} 关于\Term{conservative-functor}{保守}\Term{functor}{函子}族 $\varphi_i:E\to E_i$（$i\in I$）的概念，见第 I 讲 6.1。应当注意，当 $E,E_i$ 为 $U$-\Term{topos}{拓扑斯}，而\Term{functor}{函子} $\varphi_i$ 是\Term{correspondence}{对应}于\Term{topos-morphism}{拓扑斯态射} $f_i:E_i\to E$ 的\Term{inverse-image-functor}{逆像函子} $f_i^*$ 时，第 I 讲 6.2、6.3 和 6.4 中要求的全部\Term{exactness}{正合性}均成立；其中第 I 讲 6.2 (v) 中不带括号的情形须假设 $D$ \Term{finite}{有限}。此时，$(f_i^*)_{i\in I}$ \Term{conservative-functor}{保守}、对\Term{monomorphism}{单态射}\Term{conservative-functor}{保守}（第 I 讲 6.1）、对\Term{epimorphism}{满态射}\Term{conservative-functor}{保守}、\Term{faithful-functor}{忠实}，这些条件都\Term{equivalence}{等价}。

当\Term{functor}{函子}族 $f_i^*$ \Term{conservative-functor}{保守}时，有时也说\Term{topos-morphism}{拓扑斯态射}族 $(f_i)_{i\in I}$ \Term{conservative-functor}{保守}。特别地：
\begin{SGAresult}{定义}{6.4.1}
设 $E$ 为 $U$-\Term{topos}{拓扑斯}，$(p_i:P\to E)_{i\in I}$ 为 $E$ 的一族点。若\Term{correspondence}{对应}的\Term{fiber-functor}{纤维函子}族 $F\mapsto F_{p_i}$ 从 $E$ 到 $(U\text{-}\mathrm{Ens})$ 是\Term{conservative-functor}{保守}的（6.4.0），则称该点族\Term{conservative-family}{保守}。若\Term{topos}{拓扑斯} $E$ 存在\Term{conservative-points}{保守点族}（即\Term{topos}{拓扑斯}的全体点所成的族\Term{conservative-functor}{保守}），则说 $E$ 有\Term{enough-points}{足够多的点}。
\end{SGAresult}

\textbf{6.4.2。} 指出，迄今使用的所有 $U$-\Term{topos}{拓扑斯}都有\Term{enough-points}{足够多的点}。不过，“刻意去做”时，可以构造没有\Term{enough-points}{足够多的点}的\Term{topos}{拓扑斯}（7.2.6 e) 及 7.4）；这样的\Term{topos}{拓扑斯}在 2.2 的几何意义下必然非“空”。有\Term{enough-points}{足够多的点}的\Term{topos}{拓扑斯}必有一个以某个 $I\in U$ 为指标的\Term{conservative-points}{保守点族}（6.5）。但须注意，若 $E$ 为 $U$-\Term{topos}{拓扑斯}，则 $E$ 的点的\Term{isomorphism}{同构}类之集未必是 $U$-小的（7.3）；不过在实践中遇到的许多情形中确实如此。最后，关于\Term{enough-points}{足够多的点}的一个有趣的存在性定理（归于 P. Deligne），\Term{covering}{覆盖}了\Term{algebra}{代数}几何中遇到的所有情形。

\SourcePage{237}
\textbf{6.4.3。} 当 $(p_i)_{i\in I}$ 是\Term{topos}{拓扑斯} $E$ 的一个\Term{conservative-points}{保守点族}时，可以应用第 I 讲 6.2 的注记。它特别说明：$E$ 的两个\Term{arrow}{箭头} $u,v:F\rightrightarrows G$ 相等，当且仅当对每个 $i\in I$，在\Term{fiber}{纤维}上诱导的\Term{arrow}{箭头} $u_{p_i},v_{p_i}:F_{p_i}\rightrightarrows G_{p_i}$ 相等；$E$ 的\Term{arrow}{箭头} $u:F\to G$ 为\Term{monomorphism}{单态射}（或\Term{epimorphism}{满态射}），当且仅当对每个 $i\in I$，映射 $u_{p_i}:F_{p_i}\to G_{p_i}$ 也如此；$E$ 的\Term{object}{对象} $F$ 为\Term{initial-object}{始对象}（或\Term{final-object}{终对象}），当且仅当对每个 $i\in I$，$F_{p_i}$ 为空（或只有一点）；\Term{object}{对象} $F$ 为\Term{final-object}{终对象} $e_E$ 的\Term{subobject}{子对象}，当且仅当对每个 $i\in I$，$F_{p_i}$ 至多有一点。

\begin{SGAresult}{命题}{6.5}
 a) 设 $C$ 为 $U$-\Term{site}{景}，$(\varphi_i)_{i\in I}$ 为 $C$ 上的一族\Term{fiber-functor}{纤维函子}（6.3）。\Term{topos}{拓扑斯} $E=C^\sim$ 的相应点族 $(p_i)_{i\in I}$ \Term{conservative-functor}{保守}的充要条件为：对于 $C$ 中的任意族 $(X_j\to X)_{j\in J}$，若对每个 $i\in I$，相应族 $(\varphi_i(X_j)\to\varphi_i(X))_{j\in J}$ 都\Term{surjection}{满射}，则给定族 $(X_j\to X)_{j\in J}$ 为\Term{covering-family}{覆盖族}。

 b) 设 $E$ 为 $U$-\Term{topos}{拓扑斯}。若 $E$ 有\Term{enough-points}{足够多的点}（6.4.1），则 $E$ 有一个 $U$-小的\Term{conservative-points}{保守点族}。
\end{SGAresult}
\begin{proof}
断言 b) 是第 I 讲 7.7 的特例。为证明 a)，对由 $\varepsilon(X)$（$X\in\operatorname{ob}C$）构成的 $E$ 的\Term{generating-family}{生成族}应用第 I 讲 7.7。回顾（第 II 讲 4.4），族 $X_j\to X$ \Term{covering}{覆盖}，当且仅当族 $\varepsilon(X_j)\to\varepsilon(X)$ 为\Term{epimorphism}{满态射}族。若点族 $(p_i)_{i\in I}$ \Term{conservative-functor}{保守}，则由 6.4.3，这\Term{equivalence}{等价}于对每个 $i\in I$，族 $(X_j)_{p_i}\to(X)_{p_i}$ \Term{surjection}{满射}，即族 $\varphi_i(X_j)\to\varphi_i(X)$ \Term{surjection}{满射}。这证明了 a) 中的必要性。为证明充分性，应用第 I 讲 7.7 的判据，把问题归结为验证：若\Term{monomorphism}{单态射} $F\to\varepsilon(X)$ 使得对每个 $i\in I$，$F_{p_i}\to\varepsilon(X)_{p_i}$ 为\Term{isomorphism}{同构}，则该\Term{monomorphism}{单态射}为\Term{isomorphism}{同构}，或\Term{equivalence}{等价}地，为\Term{epimorphism}{满态射}。然而，可以找到一个\Term{covering}{覆盖}\Term{morphism}{态射}族 $\varepsilon(X_j)\to F$，再作细化后还可假定\Term{morphism}{态射} $\varepsilon(X_j)\to\varepsilon(X)$ 由\Term{morphism}{态射} $X_j\to X$ 诱导。由假设，对每个 $i\in I$，复合\Term{morphism}{态射} $(X_j)_{p_i}\to F_{p_i}\to(X)_{p_i}$ 构成一个\Term{surjective-family}{满射族}（因为是两个\Term{surjective-family}{满射族}的复合），即族 $\varphi_i(X_j)\to\varphi_i(X)$ \Term{surjection}{满射}。由假设，族 $X_j\to X$ 因而\Term{covering}{覆盖}，故族 $\varepsilon(X_j)\to\varepsilon(X)$ 为\Term{epimorphism}{满态射}族，从而 $F\to\varepsilon(X)$ 更是\Term{epimorphism}{满态射}。
\end{proof}

\begin{SGAresult}{推论}{6.5.1}
设 $C$ 为一个\Term{category}{范畴}。$C$ 上的\Term{topology}{拓扑}若使 $C$ 成为有足够多的\Term{fiber-functor}{纤维函子}的 $U$-\Term{site}{景}（即\Term{correspondence}{对应}\Term{topos}{拓扑斯}有\Term{enough-points}{足够多的点}），则只要知道由 $C$ 上的\Term{fiber-functor}{纤维函子}组成的 $\mathscr{H}om(C,\mathrm{Ens})$ 的\Term{full-subcategory}{满子范畴} $\Phi$，就完全知道了这个\Term{topology}{拓扑}。
\end{SGAresult}
\begin{proof}
事实上，由 6.5 a)，此时可以用\Term{fiber-functor}{纤维函子}族来描述\Term{covering-family}{覆盖族}。
\end{proof}
\SourcePage{238}
实际上，下面将看到（6.8.3），\Term{subcategory}{子范畴} $\Phi$ 包含于 $C$ 上的\Term{pro-representable-functor}{pro-可表示}\Term{functor}{函子}\Term{category}{范畴}中，故 $\Phi^\circ\simeq\operatorname{Point}(C^\sim)$ 在\Term{equivalence}{等价}意义下等同于 $\operatorname{Pro}\text{-}C$ 的一个\Term{full-subcategory}{满子范畴}。

可将 6.5.1 与 Giraud 定理第 II 讲 5.5 比较：该定理在 $C$ 上的\Term{topology}{拓扑}之集与 $\mathscr{H}om(C^\circ,\mathrm{Ens})$ 的某个\Term{full-subcategory}{满子范畴}之集之间建立一一\Term{correspondence}{对应}。

\begin{SGAresult}{习题}{6.5.2}
设 $C$ 为\Term{equivalence}{等价}于某个属于 $U$ 的\Term{category}{范畴}的\Term{category}{范畴}，$P$ 为 $\operatorname{Pro}(C)$ 的\Term{subcategory}{子范畴}。对每个 $p\in P$，以 $X\mapsto X_p$ 表示由 $p$ \Term{pro-representable-functor}{pro-表出}的\Term{functor}{函子}。若对每个 $p\in P$，族 $(X_{ip}\to X_p)_{i\in I}$ 都\Term{surjection}{满射}，就称 $C$ 中的族 $(X_i\to X)_{i\in I}$ 为 $P$-\Term{covering-family}{覆盖族}。证明，$C$ 上存在\Term{topology}{拓扑} $T_P$，其\Term{covering-family}{覆盖族}恰为 $P$-\Term{covering-family}{覆盖族}，且 $T_P$ 是使\Term{functor}{函子} $X\mapsto X_p$（$p\in P$）为\Term{fiber-functor}{纤维函子}的 $C$ 上最细\Term{topology}{拓扑}。证明，\Term{topology}{拓扑} $T_P$ 使 $C$ 成为有足够多的\Term{fiber-functor}{纤维函子}的\Term{site}{景}。证明，对 $C$ 上的任意\Term{topology}{拓扑} $T$，比 $T$ 更细且具有足够多的\Term{fiber-functor}{纤维函子}的\Term{topology}{拓扑} $T'$ 中有一个最粗的：它就是\Term{topology}{拓扑} $T_P$，其中 $P$ 为下述 (6.8.3) 所描述的、\Term{equivalence}{等价}于 $\operatorname{Point}(C)$ 的 $\operatorname{Pro}(C)$ 的\Term{strictly-full-subcategory}{严格满子范畴}。
\end{SGAresult}

\begin{SGAresult}{问题}{6.5.3}
刻画 $\operatorname{Pro}(C)$ 的哪些\Term{subcategory}{子范畴} $P$（严格满、在 $U$-\Term{filtered-category}{滤过}\Term{projective-limit}{投射极限}下\Term{stable}{稳定}……）可以由 $C$ 上的某个\Term{topology}{拓扑}，作为 $\operatorname{Point}(C^\sim)$ 的\Term{essential-image}{本质像}\Term{derivation}{导出}（6.8.5）。\footnote{必要条件见 9.1.8 a) 和 e)；在 $\operatorname{Top}(X)$、$X$ 为\Term{locally-noetherian}{局部诺特}的\Term{sober-space}{清醒空间}的情形，充要条件见 9.1.12 b)。}
\end{SGAresult}

\textbf{6.6。} 设 $f:E\to E'$ 为 $U$-\Term{topos-morphism}{拓扑斯态射}，由此得到\Term{canonical}{典范}\Term{functor}{函子}
\[\operatorname{Point}(f)=(p\mapsto f\circ p):\operatorname{Point}(E)\longrightarrow\operatorname{Point}(E').\]
当有两个可复合的\Term{topos-morphism}{拓扑斯态射} $E\xrightarrow fE'\xrightarrow gE'$ 时，显然有
\[\operatorname{Point}(gf)=\operatorname{Point}(g)\circ\operatorname{Point}(f):\operatorname{Point}(E)\longrightarrow\operatorname{Point}(E')\longrightarrow\operatorname{Point}(E'').\]
这就明确说明了 $\operatorname{Point}(E)$ 对\Term{topos}{拓扑斯} $E$ 的\Term{functor}{函子}性依赖（这一次并未简化用语）：若 $V$ 为满足 $U\in V$ 的\Term{universe}{宇宙}，则得到一个从属于 $V$ 的 $U$-\Term{topos}{拓扑斯}的\Term{category}{范畴}到属于 $V$ 的\Term{category}{范畴}所成\Term{category}{范畴}的（真正的！）\Term{functor}{函子}
\[\operatorname{Point}:(V\text{-}U\text{-}\mathrm{top})\longrightarrow(V\text{-}\mathrm{cat}).\]

\SourcePage{239}
\textbf{6.7。} 设 $E$ 为 $U$-\Term{topos}{拓扑斯}。$E$ 的任意\Term{object}{对象} $F$ 定义一个\Term{contravariant-functor}{反变函子}
\[p\longmapsto F_p=p^*(F):\operatorname{Point}(E)^\circ\longrightarrow(U\text{-}\mathrm{Ens}),\]
从而随 $F$ 变化得到\Term{canonical}{典范}\Term{functor}{函子}
\[E\longrightarrow\widehat{\operatorname{Point}(E)}=\mathscr{H}om(\operatorname{Point}(E)^\circ,(U\text{-}\mathrm{Ens})),\tag{6.7.1}\]
由定义（6.5），它\Term{faithful-functor}{忠实}当且仅当 $E$ 有\Term{enough-points}{足够多的点}。（\Term{functor}{函子} (6.7.1) 在形式上与 Fourier 变换有某种相似……）

设 $X$ 为 $E$ 的一个\Term{object}{对象}，因而在 $C=\operatorname{Point}(E)$ 上定义一个\Term{presheaf}{预层} $\widehat X$。于是可以定义\Term{category}{范畴}
\[C_{/\widehat X}=\operatorname{Point}(E)_{/\widehat X},\]
其\Term{object}{对象}是 $C$ 上\Term{presheaf}{预层}的\Term{category}{范畴} $\widehat C$ 中的\Term{morphism}{态射} $p\to\widehat X$，即二元组 $(p,\xi)$，其中 $p$ 为 $E$ 的点，$\xi$ 是 $\widehat X(p)=X_p$（$X$ 在 $p$ 处的\Term{fiber}{纤维}）的一个元素。现在断言，有一个\Term{canonical}{典范}的\Term{equivalence-of-categories}{范畴等价}
\[C'=\operatorname{Point}(E_{/X})\xrightarrow{\approx}C_{/\widehat X},\qquad C=\operatorname{Point}(E),\tag{6.7.2}\]
它与包含\Term{functor}{函子} $C_{/\widehat X}\to C$ 的复合就是\Term{functor}{函子}
\[\operatorname{Point}(j_X):C'=\operatorname{Point}(E_{/X})\longrightarrow C=\operatorname{Point}(E),\]
后者由\Term{localization-morphism}{局部化态射} $j_X:E_{/X}\to E$（5.2.1）的\Term{functor}{函子}性（6.6）\Term{derivation}{导出}。这只是令 $E'=P$ 所得的 5.12 的特例。

\begin{SGAresult}{推论}{6.7.3}
设 $E$ 为\Term{topos}{拓扑斯}。若 $E$ 有\Term{enough-points}{足够多的点}，则任意\Term{induced-topos}{诱导拓扑斯} $E_{/X}$（$X\in\operatorname{ob}E$）也如此。反之，若 $(X_i)_{i\in I}$ 为\Term{covering}{覆盖}\Term{final-object}{终对象}的一族 $E$ 的\Term{object}{对象}，使得对每个 $i\in I$，$E_{/X_i}$ 有\Term{enough-points}{足够多的点}，则 $E$ 有\Term{enough-points}{足够多的点}。
\end{SGAresult}
\begin{proof}
对于第一断言，设 $f:X'\to X''$ 为 $E_{/X}$ 中不是\Term{isomorphism}{同构}的\Term{morphism}{态射}，证明存在 $E_{/X}$ 的点 $q$，使 $f_q$ 不是\Term{isomorphism}{同构}。因为 $E$ 有\Term{enough-points}{足够多的点}，存在 $E$ 的点 $p$，使 $f_p:X'_p\to X''_p$ 不是\Term{isomorphism}{同构}。于是存在 $q\in X_p$，使 $f_p$ 在 $X'_p,X''_p$ 于 $X_p$ 的点 $q$ 处的\Term{fiber}{纤维} $X'_q,X''_q$ 上诱导的\Term{morphism}{态射}不是\Term{isomorphism}{同构}。但由\Term{equivalence}{等价} (6.7.2)，$q$ 可以等同于 $E_{/X}$ 的点，且刚考虑的映射 $X'_q\to X''_q$ 正是 $f:X'\to X''$ 在 $q$ 处\Term{fiber}{纤维}上诱导的映射，故得结论。

反之，假设 $X_i$ \Term{covering}{覆盖} $E$ 的\Term{final-object}{终对象}，且 $E_{/X_i}$ 有\Term{enough-points}{足够多的点}，证明 $E$ 也如此。设 $f:X'\to X''$ 是 $E$ 中不是\Term{isomorphism}{同构}的\Term{morphism}{态射}。于是存在 $i\in I$，使 $f_i:X'_i\to X''_i$
\SourcePage{240}
不是\Term{isomorphism}{同构}，其中下标 $i$ 表示\Term{restriction}{限制}到 $X_i$。于是存在 $E_{/X_i}$ 的点 $p_i$，使 $f_{ip_i}:X'_{ip_i}\to X''_{ip_i}$ 不是\Term{isomorphism}{同构}。以 $p$ 表示 $p_i$ 通过\Term{localization-morphism}{局部化态射} $E_{/X_i}\to E$ 在 $E$ 中的像，这意味着 $X'_p\to X''_p$ 不是\Term{isomorphism}{同构}，因而证明了 $E$ 有\Term{enough-points}{足够多的点}。
\end{proof}

\textbf{注记 6.7.4。} 上述论证说明，若 $(p_\alpha)$ 为 $E$ 的\Term{conservative-points}{保守点族}，则 $E_{/X}$ 中位于某个 $p_\alpha$ 之上的点所成的族（以各 $E_{p_\alpha}$ 的\Term{disjoint-union}{不交并}为指标的族）是\Term{conservative-functor}{保守}的。同样，若 $X_i$ \Term{covering}{覆盖} $E$ 的\Term{final-object}{终对象}，且对每个 $i\in I$，$P_i$ 是 $E_{/X_i}$ 的\Term{conservative-functor}{保守}点集，则 $p\in P_i$ 在\Term{localization-morphism}{局部化态射} $E_{/X_i}\to E$ 下的像所组成的 $E$ 的点集（$i\in I$ 与 $p$ 均变化）是\Term{conservative-functor}{保守}的。

\textbf{6.8。} 设 $E$ 为\Term{topos}{拓扑斯}，$p$ 为 $E$ 的点。\Term{topos}{拓扑斯} $E$ 的点 $p$ 的\Term{neighborhood}{邻域}是指二元组 $(X,u)$，其中 $X\in\operatorname{ob}E$，$u\in X_p$。由 (6.7.2)，也可以把 $u$ 解释为 $p$ 到\Term{induced-topos}{诱导拓扑斯} $E_{/X}$ 的一个点的\Term{lifting}{提升}。引入与点 $p$ \Term{correspondence}{对应}的\Term{fiber-functor}{纤维函子} $\varphi_p$，也可以把点 $p$ 的\Term{neighborhood}{邻域}解释为\Term{category}{范畴} $E_{/\varphi_p}$ 的\Term{object}{对象}（它是 $(E^\circ)^{\wedge}_{/\varphi_p}$ 中由源属于 $E\subset\widehat E$ 的\Term{object}{对象}组成的\Term{full-subcategory}{满子范畴}）。后一个\Term{category}{范畴}，即上述二元组 $(X,u)$ 的\Term{category}{范畴}，自然称为\Term{topos}{拓扑斯} $E$ 的点 $p$ 的\Term{neighborhood}{邻域}\Term{category}{范畴}，也记为 $\operatorname{Vois}(p)$。因为 $\varphi_p$ \Term{left-exact}{左正合}，且 $E$ 中\Term{finite}{有限}\Term{projective-limit}{投射极限}\Term{representable}{可表}，所以\Term{category}{范畴} $\operatorname{Vois}(p)$ 中\Term{finite}{有限}\Term{projective-limit}{投射极限}\Term{representable}{可表}，\Term{canonical}{典范}\Term{functor}{函子} $\operatorname{Vois}(p)\to E$ 与之交换（第 I 讲）。因此\Term{opposite-category}{反范畴} $\operatorname{Vois}(p)^\circ$ 是\Term{filtered-category}{滤过}的。此外，显然（第 I 讲 3.4）有关于 $F\in\operatorname{ob}E$ \Term{functor}{函子}性的\Term{canonical}{典范}\Term{isomorphism}{同构}
\[\varphi_p(F)=F_p\simeq\varinjlim_{\operatorname{Vois}(p)^\circ}F(X).\tag{6.8.1}\]
换言之，\Term{functor}{函子} $\varphi_p$“在点 $p$ 处取\Term{fiber}{纤维}”\Term{isomorphism}{同构}于由 $E$ 中点 $p$ 的\Term{neighborhood}{邻域}\Term{represented}{表出}的\Term{functor}{函子}的\Term{filtered-category}{滤过}\Term{inductive-limit}{归纳极限}。甚至可以看出，\Term{fiber-functor}{纤维函子}是\Term{ind-representable-functor}{ind-可表示}的（第 I 讲 8），也就是说，它可由 $E$ 的一个\Term{pro-object}{pro-对象} $(X_i)_{i\in I}$ \Term{represented}{表出}，其中 $I$ 为\Term{filtered-category}{滤过}\Term{ordered-set}{有序集}且 $I\in U$。事实上，由上述出处，这\Term{equivalence}{等价}于说\Term{category}{范畴} $\operatorname{Vois}(p)$ 有一个小的\Term{cofinal}{共尾}\Term{full-subcategory}{满子范畴}。设 $C$ 为 $E$ 的一个小的生成\Term{full-subcategory}{满子范畴}；断言 $\operatorname{Vois}(p)=E_{/\varphi_p}$ 中由满足 $X\in\operatorname{ob}C$ 的 $p$ 的\Term{neighborhood}{邻域} $(X,u)$ 所成的\Term{full-subcategory}{满子范畴} $C_{/\varphi_p}$（显然它是小的），在 $\operatorname{Vois}(p)$ 中\Term{cofinal}{共尾}。事实上，设 $(X,u)$ 为 $p$ 的\Term{neighborhood}{邻域}，需找到它上方的一个\Term{neighborhood}{邻域} $(X',u')$，满足 $X'\in\operatorname{ob}C$。然而，存在\Term{covering-family}{覆盖族} $X'_i\to X$，其中 $X'_i$ 属于 $C$，从而各 $X'_{ip}$ \Term{covering}{覆盖} $X_p$，故存在 $X'=X'_{i_0}$ 及 $u'\in X'_p$，使 $(X',u')$ 位于 $(X,p)$ 之上，正如所断言的。

\SourcePage{241}
\textbf{6.8.2。} 设 $C$ 为 $U$-\Term{site}{景}，$p$ 为 $C$ 的点，即\Term{topos}{拓扑斯} $C^\sim$ 的点。仍简写 $\varphi_p$，表示\Term{fiber-functor}{纤维函子} $\varphi_p$ 到 $C$ 的“\Term{restriction}{限制}”，即复合 $\varphi_p\circ\varepsilon$；对 $C$ 的\Term{object}{对象} $X$，也常写作
\[X_p=\varphi_p(X)=\varepsilon(X)_p.\]
\Term{site}{景} $C$ 中点 $p$ 的\Term{neighborhood}{邻域}是指二元组 $(X,u)$，其中 $X\in\operatorname{ob}C$，$u\in X_p=\varphi_p(X)$。这些\Term{neighborhood}{邻域}仍组成\Term{category}{范畴} $C_{/\varphi_p}$，可记为 $\operatorname{Vois}_C(p)$。证明\Term{category}{范畴} $\operatorname{Vois}_C(p)^\circ$ 仍为\Term{filtered-category}{滤过}\Term{category}{范畴}，它有一个 $U$-小的\Term{cofinal}{共尾}\Term{full-subcategory}{满子范畴}，并且有关于 $C$ 上的\Term{sheaf}{层} $F$ \Term{functor}{函子}性的\Term{isomorphism}{同构}
\[F_p\simeq\varinjlim_{\operatorname{Vois}_C(p)^\circ}F(X).\tag{6.8.3}\]
首先，重复 (6.8.1) 的论证，可见若 $C'$ 是 $C$ 的一个在\Term{topology}{拓扑}意义下生成的\Term{full-subcategory}{满子范畴}（第 II 讲 3.0.1），则 $\operatorname{Vois}_{C'}(p)^\circ$ 在 $\operatorname{Vois}_C(p)^\circ$ 中\Term{cofinal}{共尾}。取 $C'$ 为 $U$-小的，便可为证明断言而归结到 $C\in U$ 的情形。此时 $\widehat C$ 是 $U$-\Term{topos}{拓扑斯}，且有\Term{associated-sheaf}{伴随层}\Term{functor}{函子} $\underline{\phantom P}:\widehat C\to C^\sim$，由此在 $\widehat C$ 上构造复合\Term{fiber-functor}{纤维函子} $\varphi_p\underline{\phantom P}$。把 (6.8.1) 应用于\Term{topos}{拓扑斯} $\widehat C$ 及其生成\Term{full-subcategory}{满子范畴} $C$，可见\Term{category}{范畴} $C^\circ_{/\varphi_p}$（它显然\Term{isomorphism}{同构}于 $\operatorname{Vois}_C(p)^\circ$）是\Term{filtered-category}{滤过}的，且有关于 $C$ 上\Term{presheaf}{预层} $P$ \Term{functor}{函子}性的\Term{isomorphism}{同构}
\[\varphi_p(\underline P)\simeq\varinjlim_X P(X).\]
若 $F$ 是 $C$ 上的\Term{sheaf}{层}，则有 $\underline F=F$，将上述\Term{isomorphism}{同构}应用于作为\Term{presheaf}{预层}的 $F$，即得 (6.8.3)。同时，这一论证还给出更一般的公式
\[(\underline P)_p\simeq\varinjlim_{\operatorname{Vois}_C(p)^\circ}P(X),\tag{6.8.4}\]
即关于 $C$ 上任意\Term{presheaf}{预层} $P$ 的\Term{functor}{函子}性\Term{isomorphism}{同构}，至少当 $C$ 小时如此。一般情形可通过引入满足 $C\in V$ 的\Term{universe}{宇宙} $V$，把 $C$ 上的 $\varphi_p$ 看成取值于 $(V\text{-}\mathrm{Ens})$ 的\Term{fiber-functor}{纤维函子}，并利用\Term{adjunction}{伴随}\Term{presheaf}{预层}的形成与扩大\Term{universe}{宇宙}的相容性（第 II 讲 3.6）而得出。

\textbf{6.8.5。} 已给\Term{topos}{拓扑斯} $C^\sim$ 的每个点 $p$ 配以 $U$-\Term{site}{景} $C$ 的一个\Term{pro-object}{pro-对象}；它由\Term{canonical}{典范}\Term{functor}{函子} $\operatorname{Vois}_C(p)\to C$ 定义，这里利用了 $p$ 在 $C$ 中的\Term{neighborhood}{邻域}\Term{category}{范畴}是\Term{cofiltered-category}{余滤过}的且有一个小的\Term{cofinal}{共尾}\Term{full-subcategory}{满子范畴}。立即可验证，随 $p$ 变化，由此得到\Term{functor}{函子}
\[\operatorname{Point}(C^\sim)\longrightarrow\operatorname{Pro}(C).\tag{6.8.5.1}\]
此\Term{functor}{函子}\Term{fully-faithful-functor}{全忠实}。事实上，通过\Term{correspondence}{对应}的\Term{fiber-functor}{纤维函子}，它\Term{solvable}{可解}释为\Term{functor}{函子} $\operatorname{Fib}(C)\to\operatorname{Fib}(C^\sim)\to\operatorname{Pro}(C)^\circ$。第一个\Term{functor}{函子}
\SourcePage{242}
由 4.9.4 是\Term{equivalence}{等价}（在 6.3 中已对\Term{fiber-functor}{纤维函子}回顾），第二个\Term{functor}{函子}由\Term{pro-object}{pro-对象}的一般性质（第 I 讲 8）是\Term{fully-faithful-functor}{全忠实}的。

\textbf{6.8.6。} 特别地，设 $C$ 为 $U$-小且配备\Term{chaotic-topology}{混沌拓扑}，因此 $C^\sim=\widehat C$。断言此时上述\Term{functor}{函子}甚至是\Term{equivalence-of-categories}{范畴等价}
\[\operatorname{Point}(\widehat C)\xrightarrow{\approx}\operatorname{Pro}(C).\tag{6.8.6.1}\]
事实上，只需证明它\Term{essentially-surjective-functor}{本质满射}；这是因为 $\widehat C$ 上形如 $P\mapsto P(X)$ 的\Term{functor}{函子}都是\Term{fiber-functor}{纤维函子}，以及\Term{topos}{拓扑斯}上的\Term{fiber-functor}{纤维函子}的任意\Term{filtered-category}{滤过}\Term{inductive-limit}{归纳极限}显然仍为\Term{fiber-functor}{纤维函子}。

把 $C$ 等同于 $\operatorname{Pro}(C)$ 的\Term{full-subcategory}{满子范畴}（第 I 讲 8），则由 (6.8.5.2) 的一个\Term{quasi-inverse-functor}{拟逆}\Term{functor}{函子}诱导的\Term{functor}{函子} $C\to\operatorname{Point}(\widehat C)$，显然\Term{isomorphism}{同构}于 (4.6.2.2) 中已考虑过的\Term{canonical}{典范}\Term{functor}{函子}；利用 $\operatorname{Point}(\widehat C)$ 在小\Term{projective-limit}{投射极限}下\Term{stable}{稳定}这一事实，可将上述\Term{quasi-inverse-functor}{拟逆}看作该\Term{functor}{函子}的\Term{canonical-extension}{典范延拓}（第 I 讲 8）。

\textbf{6.8.7。} 一般地，设 $(X_i)_{i\in I}$ 为 $U$-\Term{site}{景} $C$ 的一个\Term{pro-object}{pro-对象}，从而在 $C^\sim$ 上有\Term{functor}{函子}
\[\varphi(F)=\varinjlim_I F(X_i).\tag{6.8.7.1}\]
鉴于 (6.8.5)，自然要问何时此\Term{functor}{函子}是\Term{fiber-functor}{纤维函子}。以下条件\Term{equivalence}{等价}：

 (i) 上述\Term{functor}{函子} $\varphi$ 是 $C^\sim$ 上的\Term{fiber-functor}{纤维函子}。

 (ii) $\varphi$ 到 $C$ 的“\Term{restriction}{限制}”是\Term{site}{景} $C$ 上的\Term{fiber-functor}{纤维函子}（6.3）。

 (iii) 对 $C$ 的\Term{object}{对象} $Y$ 的任意\Term{covering-family}{覆盖族} $Y_\alpha\to Y$、任意 $i_0\in I$ 及任意\Term{morphism}{态射} $X_{i_0}\to Y$，都存在 $i\geq i_0$、一个 $\alpha$ 和一个\Term{morphism}{态射} $X_i\to Y_\alpha$，使图
\[\begin{tikzcd}X_i\arrow[r]\arrow[d]&Y_\alpha\arrow[d]\\X_{i_0}\arrow[r]&Y\end{tikzcd}\]
交换。

事实上，条件 (iii) 正表达了 $\varphi|C$ 把\Term{covering-family}{覆盖族}变为\Term{covering-family}{覆盖族}。蕴含 (i)$\Rightarrow$(ii)$\Rightarrow$(iii) 显然，故只需证明 (iii)$\Rightarrow$(i)。因为\Term{functor}{函子} $\varphi$ 显然\Term{left-exact}{左正合}，只需证明（4.9.4）它把 $C^\sim$ 中的\Term{covering-family}{覆盖族} $F_\alpha\to F$ 变为\Term{covering-family}{覆盖族}，即对任意 $i_0$ 及 $x_0\in F(X_{i_0})$，存在 $i\geq i_0$、一个 $\alpha$ 及 $y\in F_\alpha(X_i)$，使 $x_0$ 和 $y$ 在 $F(X_i)$ 中的像相同。
\SourcePage{243}
因为 $F_\alpha\to F$ \Term{covering}{覆盖}，存在 $X_{i_0}$ 的\Term{covering-family}{覆盖族} $Z_\gamma\to X_{i_0}$，使对每个 $\gamma$，$x_0$ 在 $F(Z_\gamma)$ 中的\Term{inverse-image}{逆像}可以\Term{lifting}{提升}为某个 $F_\alpha(Z_\gamma)$ 的元素。由假设 (iii)，存在 $i\geq i_0$ 和 $X_{i_0}$-\Term{morphism}{态射} $X_i\to Z_\gamma$。于是 $x_0$ 在 $F(X_i)$ 中的像可以\Term{lifting}{提升}为某个 $F_\alpha(X_i)$ 的元素 $y$，证毕。

\begin{SGAresult}{习题}{6.9}
设 $C$ 为属于 $U$ 的\Term{site}{景}，$\varphi$ 为 $C$ 上的\Term{fiber-functor}{纤维函子}，由 $\operatorname{Pro}(C)$ 的\Term{object}{对象} $(X_i)_{i\in I}$ \Term{pro-representable-functor}{pro-表出}。对于任意指标 $i\in I$、$C$ 的任意\Term{object}{对象} $Y$、任意\Term{morphism}{态射} $u:X_i\to Y$ 及 $Y$ 的任意\Term{covering}{覆盖}筛 $R$，选择指标 $j\geq i$，使复合 $X_j\to X_i\xrightarrow uY$ 经 $R$ \Term{resolution}{分解}。对固定的 $i$，令 $J(i)$ 为由 $i$ 以及随 $Y,u,R$ 变化所得的 $j$ 所组成的\Term{set}{集合}；对 $I$ 的\Term{subset}{子集} $I_0$，同样令 $J(I_0)$ 为 $i\in I_0$ 时各 $J(i)$ 的并。另一方面，若 $I_0$ 为 $I$ 的\Term{finite}{有限}\Term{subset}{子集}，令 $m(I_0)\in I$ 为 $I_0$ 的一个上界；对 $I$ 的任意\Term{subset}{子集} $I_0$，令 $M(I_0)$ 为各 $M(I'_0)$ 的并，其中 $I'_0$ 遍历 $I_0$ 的\Term{finite}{有限}\Term{subset}{子集}之集。假定函数 $I_0\mapsto m(I_0)$ 的选择满足：当 $I_0$ 只有一个元素 $i$ 时，$m(I_0)=i$；于是对 $I$ 的任意\Term{subset}{子集} $I_0$，均有 $I_0\subset M(I_0)$。考虑从 $I$ 的幂集到自身的映射 $MJ$ 的迭代 $(MJ)^n$；对 $I$ 的任意\Term{subset}{子集} $I_0$，令 $P(I_0)$ 为各 $(MJ)^n(I_0)$ 的并。

 a) 证明，对 $I$ 的任意\Term{subset}{子集} $I_0$，$P(I_0)$ 关于诱导次序是 $I$ 的\Term{filtered-category}{滤过}\Term{subset}{子集}；且当 $I_0$ 遍历 $I$ 的\Term{finite}{有限}\Term{subset}{子集}之集时，$I$ 是 $P(I_0)$ 的递增\Term{filtered-category}{滤过}并。

 b) 证明，存在只依赖于\Term{cardinal}{基数} $\operatorname{card}\operatorname{fl}(C)=a$ 的\Term{cardinal}{基数} $c\in U$，使对 $I$ 的任意\Term{finite}{有限}\Term{subset}{子集} $I_0$，有 $\operatorname{card}(P(I_0))\leq c$。（若 $a$ 无限，可取 $c=2^a$。）

 c) 利用 6.8.7 证明，对 $I$ 的任意\Term{subset}{子集} $I_0$，由\Term{pro-object}{pro-对象} $(X_i)_{i\in I_0}$ \Term{pro-representable-functor}{pro-表出}的 $C$ 上的\Term{functor}{函子} $\varphi_{I_0}$ 为 $C$ 上的\Term{fiber-functor}{纤维函子}；且当 $I_0$ 遍历 $I$ 的\Term{finite}{有限}\Term{subset}{子集}之集时，\Term{fiber-functor}{纤维函子} $\varphi$ 是\Term{functor}{函子} $\varphi_{I_0}$ 的\Term{filtered-category}{滤过}\Term{inductive-limit}{归纳极限}。

 d) 推出，$\operatorname{Fib}(C)\simeq\operatorname{Fib}(C)$ 有一个小的\Term{full-subcategory}{满子范畴} $\Phi$，使 $\operatorname{Fib}(C)$ 的任意\Term{object}{对象}均为 $\Phi$ 的\Term{object}{对象}的\Term{filtered-category}{滤过}\Term{inductive-limit}{归纳极限}。（若 $a$ 无限，可取 $\Phi$ 的\Term{cardinal}{基数} $\leq2^{2^a}$；若 $a$ \Term{finite}{有限}，当然可以取 $\Phi=C$。）
\end{SGAresult}

\begin{SGAresult}{习题}{6.10}
设 $C$ 为 $U$-\Term{site}{景}，$D$ 为其中小的\Term{filtered-category}{滤过}\Term{inductive-limit}{归纳极限}均\Term{representable}{可表}的 $U$-\Term{category}{范畴}，$\operatorname{Fais}(C,D)$ 为 $C$ 上取值于 $F$ 的\Term{sheaf}{层}\Term{category}{范畴}（第 II 讲 6.1）。推广 (6.8.3)，定义一个“\Term{fiber-functor}{纤维函子}”
\[F\longmapsto F_p:\operatorname{Fais}(C,D)\longrightarrow D.\tag{6.10.1}\]
若 $D$ 中的\Term{finite}{有限}\Term{projective-limit}{投射极限}\Term{representable}{可表}且与\Term{filtered-category}{滤过}\Term{inductive-limit}{归纳极限}交换，则上述\Term{functor}{函子}\Term{left-exact}{左正合}。
\end{SGAresult}

\SGAsection{7}{纤维函子与拓扑斯的点的例子}
\SourcePage{244}
\textbf{7.1。\Term{topological-space}{拓扑空间} $X$ 的 $\operatorname{Top}(X)$ 的情形。} 设 $x$ 为 $X$ 的一个点。把单点集 $\{x\}$ 看成\Term{topological-space}{拓扑空间}，并考虑包含 $i_x:\{x\}\to X$。由 4.1.1，它定义一个\Term{topos-morphism}{拓扑斯态射}（在唯一\Term{isomorphism}{同构}意义下确定）
\[\operatorname{Top}(i_x):\operatorname{Top}(\{x\})\longrightarrow\operatorname{Top}(X).\tag{7.1.1}\]
另一方面，可以把 $P$ 等同于 $\operatorname{Top}(\{x\})$（2.2），于是得到 $\operatorname{Top}(X)$ 的一个点 $p_x$：
\[p_x:P\longrightarrow\operatorname{Top}(X),\tag{7.1.2}\]
它由 $x$ 在唯一\Term{isomorphism}{同构}意义下定义。\Term{correspondence}{对应}的\Term{fiber-functor}{纤维函子}由熟知的公式 [TF] 给出；该公式也立即由第 I 讲 5.1 得出：
\[p_x^*(F)=F_x=F_{p_x}=\varinjlim_{U\ni x}F(U),\tag{7.1.3}\]
其中\Term{limit}{极限}按 $x$ 在 $X$ 中的开\Term{neighborhood}{邻域} $U$ 取。

我们知道 $\operatorname{Top}(X)$ \Term{isomorphism}{同构}于 $\operatorname{Top}(X_{\mathrm{sob}})$（4.2.1），所以更一般地，$X_{\mathrm{sob}}$ 的每个点，即 $X$ 的每个\Term{irreducible}{不可约}\Term{closed-subset}{闭子集} $Z$，都定义 $\operatorname{Top}(X)$ 的一个点，或者说 $\operatorname{Top}(X)$ 上的\Term{fiber-functor}{纤维函子}，仍记为 $F\mapsto F_Z$。回到它在 $X_{\mathrm{sob}}$ 上由公式 (7.1.3) 给出的定义，得到
\[F_Z=\varinjlim_{U\cap Z\ne\varnothing}F(U),\tag{7.1.4}\]
其中\Term{inductive-limit}{归纳极限}按 $X$ 中与 $Z$ 相交的\Term{open-subset}{开集} $U$ 取。

\textbf{7.1.5。} 众所周知，\Term{fiber-functor}{纤维函子} $F\mapsto F_x$（$x\in X$）本身就组成一个\Term{conservative-functor}{保守}族。把 $X$ 上的\Term{sheaf}{层}解释为\Term{etale-space}{平展空间}时，这特别明显：此时 $F$ 在 $x$ 处的\Term{fiber}{纤维}正是 $X$ 上\Term{fibred-space}{纤维空间}意义下的\Term{fiber}{纤维}。注意，所考虑的\Term{conservative-functor}{保守}\Term{fiber-functor}{纤维函子}族的指标集为 $X$，因此是 $U$-小的。

\textbf{7.1.6。} 由 4.2.3，\Term{category}{范畴} $\operatorname{Point}(\operatorname{Top}(X))$ \Term{equivalence}{等价}于 $X_{\mathrm{sob}}$ 按\Term{specialization}{特化}关系排序所得的\Term{ordered-set}{有序集}所\Term{correspondence}{对应}的\Term{category}{范畴}。换言之，$\operatorname{Top}(X)$ 上的任意\Term{fiber-functor}{纤维函子}都\Term{isomorphism}{同构}于\Term{fiber-functor}{纤维函子} $F\mapsto F_Z$，其中 $Z$ 是 $X_{\mathrm{sob}}$ 的唯一确定元素，即 $X$ 的唯一确定的\Term{irreducible}{不可约}\Term{closed-subset}{闭子集}；另一方面，若 $Z,Z'\in X_{\mathrm{sob}}$，则\Term{set}{集合} $\operatorname{Hom}(F_Z,F_{Z'})$ 为空或只有一个元素，后一情形当且仅当 $Z$ 为 $X_{\mathrm{sob}}$ 中 $Z'$ 的\Term{specialization}{特化}，即当且仅当作为 $X$ 的\Term{subset}{子集}有 $Z\subset Z'$。若 $X$ 本身清醒，则这些陈述中可把 $X_{\mathrm{sob}}$ 换成 $X$ 本身。

注意，$\operatorname{Top}(X)$ 的\Term{fiber-functor}{纤维函子}的\Term{automorphism-group}{自同构群}（与 $\operatorname{Top}(X)$ 的\Term{correspondence}{对应}点的\Term{automorphism-group}{自同构群}相反）总是
\SourcePage{245}
\Term{unit-group}{单位群}（更一般地，由 4.2.3，\Term{topological-space}{拓扑空间}所\Term{correspondence}{对应}的\Term{topos}{拓扑斯}之间任意\Term{topos-morphism}{拓扑斯态射}的\Term{automorphism-group}{自同构群}也如此）。这是所考虑特殊情形独有的现象：见例 7.2 及第 VIII 讲 7。对于“模问题”所\Term{correspondence}{对应}的\Term{topos}{拓扑斯}（例如见 [10]），\Term{fiber-functor}{纤维函子}的\Term{automorphism-group}{自同构群}有一个很好的解释，即所要分类的（\Term{algebraically-closed-field}{代数闭域}上的）\Term{algebraic-structure}{代数结构}的\Term{automorphism-group}{自同构群}。

\textbf{7.1.7。} 刚才已看到，如何借助\Term{topos}{拓扑斯} $\operatorname{Top}(X)$（甚至只需知道它的\Term{equivalence-class}{等价类}），把\Term{sober-space}{清醒空间} $X$（或与任意\Term{topological-space}{拓扑空间}\Term{correspondence}{对应}的\Term{sober-space}{清醒空间} $X_{\mathrm{sob}}$）至少作为按\Term{specialization}{特化}关系排序的\Term{ordered-set}{有序集}，重构为 $\operatorname{Top}(X)$ 的“点”的\Term{isomorphism}{同构}类之集。由 2.1 中所述，还可如下重构 $X$ 的\Term{topology}{拓扑}，即其\Term{open-subset}{开集}族：对 $E$ 的\Term{final-object}{终对象} $e_E$ 的每个\Term{subobject}{子对象} $U$，令 $\operatorname{Pt}(U)$ 为满足 $U_x\ne\varnothing$ 的 $x\in X$ 之集（由 (6.7.2)，它也等同于\Term{induced-topos}{诱导拓扑斯} $E_{/U}$ 的点的\Term{isomorphism}{同构}类之集）。于是 $U\mapsto\operatorname{Pt}(U)$ 是从 $e_E$ 的\Term{subobject}{子对象}之集到 $X$ 的\Term{open-subset}{开集}之集的\Term{ordered-set}{有序集}\Term{isomorphism}{同构}。

\textbf{7.1.8。} 7.1.6 对\Term{topological-space}{拓扑空间} $X$ 所定义的\Term{topos}{拓扑斯} $\operatorname{Top}(X)$ 的点\Term{category}{范畴}的确定，促使我们对任意\Term{topos}{拓扑斯} $E$ 的点 $p,p'$ 采用如下术语：若存在从 $p'$ 到 $p$ 的\Term{morphism}{态射}（按 6.1 的\Term{category}{范畴} $\operatorname{Point}(E)$ 的意义），则称 $p$ 为 $p'$ 的\Term{specialization}{特化}，或称 $p'$ 为 $p$ 的\Term{generalization}{泛化}。这些关系仍具\Term{transitivity}{传递性}，但利用 6.8.6 很容易找到 $p,p'$ 互为\Term{specialization}{特化}而不\Term{isomorphism}{同构}（当然更不相等）的例子。\Term{topos}{拓扑斯} $E$ 的点 $p$ 的\Term{generalization}{泛化}\Term{category}{范畴}，即\Term{category}{范畴} $\operatorname{Point}(E)_{/p}$，在某些方面起着与在\Term{scheme}{概形}的某点处作\Term{localization}{局部化}相似的\Term{action}{作用}。第 VI 章将通过构造 $E$ 在点 $p$ 处的\Term{localization}{局部化}\Term{topos}{拓扑斯}，使这一\Term{action}{作用}更加明确；其点\Term{category}{范畴}\Term{canonical}{典范}\Term{equivalence}{等价}于 $p$ 的\Term{generalization}{泛化}\Term{category}{范畴}。

\begin{SGAresult}{习题}{7.1.9}
设 $E$ 为 $U$-\Term{topos}{拓扑斯}。证明，$E$ \Term{equivalence}{等价}于某个形如 $\operatorname{Top}(X)$ 的\Term{topos}{拓扑斯}（其中 $X$ 是属于 $U$ 的\Term{topological-space}{拓扑空间}），当且仅当它满足下述两个条件：

 a) \Term{final-object}{终对象} $e_E$ 的\Term{subobject}{子对象}族是\Term{topos}{拓扑斯} $E$ 的\Term{generating-family}{生成族}。

 b) $E$ 有\Term{enough-points}{足够多的点}（6.5）。（此条件并非多余：见 7.4。）

与 7.8 c) 比较。
\end{SGAresult}

\begin{SGAresult}{习题}{7.1.10}
设 $X$ 为带有算\Term{subgroup}{子群} $G$ 的\Term{topological-space}{拓扑空间}（$X,G\in U$），并令 $E=\operatorname{Top}(X,G)$（2.3）。
\SourcePage{246}
 a) 证明，对于 $E$ 的每个\Term{object}{对象} $(X',G)$，\Term{induced-topos}{诱导拓扑斯} $E_{/(X',G)}$ \Term{canonical}{典范}\Term{equivalence}{等价}于\Term{topos}{拓扑斯} $\operatorname{Top}(X',G)$（比较 5.7）。

 b) 当 $G$ 在 $X'$ 上的\Term{action}{作用}\Term{proper}{固有}且自由时（Bourbaki，《一般\Term{topology}{拓扑}》第 III 章 4），证明带\Term{operator}{算子}\Term{space}{空间}的\Term{morphism}{态射} $(X',G)\to(X'/G,e)$ 诱导（4.1.2）\Term{topos}{拓扑斯}\Term{equivalence}{等价}
\[\operatorname{Top}(X',G)\longrightarrow\operatorname{Top}(X'/G).\]
 c) 由 b) 推出，存在 $E$ 的\Term{object}{对象} $Z$，使\Term{induced-topos}{诱导拓扑斯} $E_{/Z}$ \Term{equivalence}{等价}于 $\operatorname{Top}(X)$，且\Term{localization-functor}{局部化函子} $E\to E_{/Z}$ \Term{isomorphism}{同构}于“遗忘 $G$ 的\Term{action}{作用}”的\Term{functor}{函子}。（仿照 5.8.3 的论证。）

 d) 由 c) 推出，$E$ 上任意\Term{fiber-functor}{纤维函子}均由 $\operatorname{Top}(X)$ 上的\Term{fiber-functor}{纤维函子}通过“遗忘 $G$ 的\Term{action}{作用}”的\Term{functor}{函子}诱导，从而可由 $X_{\mathrm{sob}}$ 的一个点定义。

 e) 用带\Term{operator}{算子}\Term{space}{空间} $(X_{\mathrm{sob}},G)$ 确定\Term{category}{范畴} $\operatorname{Point}(E)$ 的结构（在\Term{equivalence}{等价}意义下）。特别地推出，$X_{\mathrm{sob}}$ 的两点定义 $E$ 上\Term{isomorphism}{同构}的\Term{fiber-functor}{纤维函子}，当且仅当它们在 $G$ 的\Term{action}{作用}下\Term{conjugacy}{共轭}。
\end{SGAresult}

\begin{SGAresult}{习题}{7.1.11}
设 $X\in U$ 为\Term{topological-space}{拓扑空间}。证明 $V$-\Term{topos}{拓扑斯} $\operatorname{TOP}(X)$（2.5）有\Term{enough-points}{足够多的点}。更准确地说，对 $\operatorname{TOP}(X)$ 的每个\Term{object}{对象} $X'$ 及每个 $x'\in X'$，定义 $\operatorname{TOP}(X)$ 上的\Term{fiber-functor}{纤维函子} $F\mapsto F_{X',x'}$，使它只依赖于 $X$ 上\Term{space}{空间} $(X',x')$ 的“芽”；并证明此\Term{fiber-functor}{纤维函子}族是\Term{conservative-functor}{保守}的（且以一个 $V$-小的指标集为指标）。利用一个适当的\Term{filtered-category}{滤过}\Term{projective-system}{投射系} $(X'_i,x'_i)_{i\in I}$，其中 $I$ 不是 $U$-小的，证明这样不能得到 $V$-\Term{topos}{拓扑斯} $\operatorname{TOP}(X)$ 的所有\Term{fiber-functor}{纤维函子}。证明，这样的\Term{fiber-functor}{纤维函子}的\Term{isomorphism}{同构}类之集的\Term{cardinal}{基数}不属于 $V$。
\end{SGAresult}

\textbf{7.2。\Term{classifying-topos}{分类拓扑斯} $B_G$ 的点。} 设 $E$ 为 $U$-\Term{topos}{拓扑斯}，$G$ 为 $E$ 中的\Term{group}{群}，从而有\Term{classifying-topos}{分类拓扑斯} $B_G$（2.4）。若 $e$ 为单点\Term{group}{群}，则\Term{morphism}{态射} $e\to G\to e$ 定义（4.5）\Term{topos-morphism}{拓扑斯态射}（这里利用 $B_e\approx E$）
\[E\longrightarrow B_G\longrightarrow E,\tag{7.2.1}\]
其复合\Term{isomorphism}{同构}于 $\operatorname{id}_E$。因此取点\Term{category}{范畴}得到\Term{functor}{函子}
\[\operatorname{Point}(E)\longrightarrow\operatorname{Point}(B_G)\longrightarrow\operatorname{Point}(E),\tag{7.2.2}\]
其复合\Term{isomorphism}{同构}于恒等\Term{functor}{函子}。利用 (7.2.1) 的第一个\Term{topos-morphism}{拓扑斯态射}等同于相对于 $B_G$ 的\Term{object}{对象} $E_G=G$ 的\Term{localization-morphism}{局部化态射}（5.8），由 (6.7.2)，(7.2.2) 的第一个\Term{functor}{函子}等同于自然的“包含”\Term{functor}{函子}
\[\operatorname{Point}(B_G)_{/\widehat E_G}\longrightarrow\operatorname{Point}(B_G),\tag{7.2.3}\]
\SourcePage{247}
其中 $\widehat E_G$ 表示 $\operatorname{Point}(B_G)$ 上的\Term{presheaf}{预层} $p'\mapsto(E_G)_{p'}$。因为 $E_G$ 显然\Term{covering}{覆盖} $B_G$ 的\Term{final-object}{终对象}，它的\Term{fiber}{纤维} $(E_G)_p$ \Term{nonempty}{非空}，因此 (7.2.3) \Term{essentially-surjective-functor}{本质满射}。这意味着，$B_G$ 上任意\Term{fiber-functor}{纤维函子}都由 $E$ 的\Term{fiber-functor}{纤维函子}通过“遗忘 $G$ 的\Term{action}{作用}”的\Term{functor}{函子} $B_G\to E$ 诱导（在\Term{isomorphism}{同构}意义下）。特别地，若 $E$ 有\Term{enough-points}{足够多的点}（或一个小的\Term{conservative-points}{保守点族}），则 $B_G$ 也如此。

\textbf{注记 7.2.4。} 还可以进一步，用\Term{category}{范畴} $C=\operatorname{Point}(E)$ 及其上的\Term{group}{群}值\Term{presheaf}{预层}
\[\widehat G:p\longmapsto G_p:C^\circ\longrightarrow(\mathrm{Groupes})\]
确定\Term{category}{范畴} $C'=\operatorname{Point}(B_G)$ 的结构（在\Term{equivalence}{等价}意义下）。事实上，定义新\Term{category}{范畴} $C_1$，使它与 $C$ 有相同的\Term{object}{对象}，并且对 $C$ 的两个\Term{object}{对象} $p,q$，有
\[\operatorname{Hom}_{C_1}(p,q)=\operatorname{Hom}_C(p,q)\times\widehat G(p),\]
而 $C_1$ 中\Term{arrow}{箭头}的复合由 $C$ 中的复合及\Term{presheaf}{预层} $\widehat G$ 的\Term{group}{群}运算诱导。给定取值于任意\Term{category}{范畴} $D$ 的\Term{functor}{函子} $\varphi_1:C_1\to D$，于是\Term{equivalence}{等价}于给定\Term{functor}{函子} $\varphi:C\to D$，并在其上配备\Term{presheaf}{预层} $\widehat G$ 的\Term{action}{作用}（即对每个 $p\in\operatorname{ob}C$，给出 $\widehat G(p)$ 在 $\varphi(p)$ 上的\Term{action}{作用}，并关于变化的 $p$ 满足一个显然的\Term{functor}{函子}性条件）。利用这一观察，可定义\Term{canonical}{典范}\Term{functor}{函子}
\[\varphi_1:C_1\longrightarrow C',\tag{7.2.4.1}\]
它\Term{correspondence}{对应}于 (7.2.2) 中的\Term{functor}{函子} $\varphi:C\to C'$：把 $E$ 的每个点 $p$ 映到 $B_G$ 上所诱导的点 $\varphi(p)$，并在后者上配备 $\widehat G(p)=G_p$ 的自然\Term{action}{作用}；当 $X$ 遍历 $B_G$ 时，这些\Term{action}{作用}来自 $G_p$ 在 $X_p$ 上的\Term{action}{作用}。留给读者证明 (7.2.4.1) 为\Term{equivalence-of-categories}{范畴等价}：利用 (7.2.2) 的第一个\Term{functor}{函子}的结构 (7.2.3)，并注意自然包含\Term{functor}{函子} $C\to C_1$ 对于 $C'$ 上的\Term{presheaf}{预层} $P_0=\widehat E_G$ 的\Term{inverse-image}{逆像} $P_1$ 有类似结构。

\textbf{7.2.5。} 特别地可知，$B_G$ 的点的\Term{isomorphism}{同构}类恰与 $E$ 的点的\Term{isomorphism}{同构}类\Term{correspondence}{对应}：\Term{classifying-topos}{分类拓扑斯} $E_G$ 的每个点都由 $E$ 的一个点诱导，所差\Term{isomorphism}{同构}不唯一。特别地，当 $E$ 为\Term{punctual-topos}{点拓扑斯} $P$，从而 $G$ 为普通\Term{group}{群}时，\Term{category}{范畴} $\operatorname{Point}(B_G)$ 为\Term{fundamental-group}{基本群}是 $G$ 的\Term{connected}{连通}\Term{groupoid}{群胚}：$B_G$ 上任意\Term{fiber-functor}{纤维函子}都（非\Term{canonical}{典范}地）\Term{isomorphism}{同构}于“遗忘 $G$ 的\Term{action}{作用}”的\Term{functor}{函子} $\omega$，而此\Term{functor}{函子}的\Term{endomorphism}{自同态}\Term{monoid}{幺半群}为\Term{group}{群} $G$（所以 $\omega$ 的每个\Term{endomorphism}{自同态}均为自\Term{isomorphism}{同构}）。

\SourcePage{248}
\begin{SGAresult}{习题}{7.2.6}
 a) 设 $(\mathrm{Tors})$ 为二元组 $(\Gamma,P)\in U$ 的\Term{category}{范畴}，其中 $\Gamma$ 为\Term{group}{群}，$P$ 为 $\Gamma$ 下的右\Term{torsor}{挠子}。\Term{functor}{函子} $(\Gamma,P)\mapsto\Gamma$
\[(\mathrm{Tors})\longrightarrow(\mathrm{Groupes})\tag{7.2.6.1}\]
为余\Term{fiber}{纤维}化\Term{functor}{函子}。沿用 7.2.4 的记号，考虑\Term{functor}{函子} $\widehat G:C^\circ\to(\mathrm{Groupes})$；令 $D'$ 为在 $C^\circ$ 上作为\Term{cofibred-category}{余纤维范畴} (7.2.6.1) 的\Term{inverse-image}{逆像}的\Term{cofibred-category}{余纤维范畴}，$D$ 为 $C$ 上\Term{correspondence}{对应}的\Term{fibred-category}{纤维范畴}，它与 $D'$ 的各\Term{fibred-category}{纤维范畴}相同，所以对 $p\in\operatorname{ob}C$，有
\[D_p=\text{$G_p$ 下的右挠子的范畴}.\]
构造互为\Term{quasi-inverse-functor}{拟逆}的\Term{canonical}{典范}\Term{functor}{函子}
\[D\longrightarrow C'=\operatorname{Point}(B_G),\qquad C'\longrightarrow D.\tag{7.2.6.2}\]
特别地，推出 $B_G$ 中位于 $E$ 的给定点 $p$ 上方的点所成\Term{category}{范畴}，\Term{canonical}{典范}\Term{equivalence}{等价}于\Term{group}{群} $G_p$ 下的右\Term{torsor}{挠子}\Term{category}{范畴}。

 b) 设 $\mathscr G=(G_i)_{i\in I}$ 为 $U$-\Term{topos}{拓扑斯} $E$ 中\Term{group}{群}的一个严格\Term{projective-system}{投射系}（$I\in U$ 为\Term{filtered-category}{滤过}\Term{preordered-set}{预序集}）；仿照 2.7.1 的定义得到\Term{classifying-topos}{分类拓扑斯} $B_{\mathcal U}$。仿照 a) 确定 $B_{\mathscr G}$ 的点\Term{category}{范畴}。（注意：把\Term{category}{范畴} $(\mathrm{Tors})$ 和 $(\mathrm{Groupes})$ 换成这些\Term{category}{范畴}中以 $I$ 为指标的\Term{projective-system}{投射系}\Term{category}{范畴}。）推出，若 $I$ 有可数\Term{cofinal}{共尾}集，则 $B_{\mathscr G}$ 上任意\Term{fiber-functor}{纤维函子}都\Term{isomorphism}{同构}于从 $E$ 的\Term{fiber-functor}{纤维函子}诱导的\Term{functor}{函子}（通过“遗忘 $G$ 的\Term{action}{作用}”的\Term{functor}{函子}）；后者在不唯一的\Term{isomorphism}{同构}意义下确定。

 c) 取 $E$ 为\Term{punctual-topos}{点拓扑斯}，从而 $G$ 为普通\Term{pro-group}{pro-群}；证明，若 $I$ 任意，但 $\mathscr G$ 为\Term{profinite}{投有限}的（即各 $G_i$ \Term{finite}{有限}），则 b) 的结论仍成立：每个\Term{fiber-functor}{纤维函子}均\Term{isomorphism}{同构}于“遗忘 $G$ 的\Term{action}{作用}”的\Term{functor}{函子}。

 d) 仍取 $E$ 为\Term{punctual-topos}{点拓扑斯}。对适当的\Term{projective-system}{投射系} $\mathscr G=(G_i)_{i\in I}$，给出 $B_{\mathscr G}$ 上一个不\Term{isomorphism}{同构}于“遗忘 $G$ 的\Term{action}{作用}”的\Term{functor}{函子}的\Term{fiber-functor}{纤维函子}的例子。

 e) 把 $\mathscr G=(G_i)$ 看成\Term{topos}{拓扑斯} $\widehat I$ 中的\Term{group}{群}，设 $\mathscr T$ 为 $I$ 上以 $G$ 为带（lien）的\Term{gerbe}{凝集} [3]。说明如何用\Term{gerbe}{凝集} $T$“扭曲”\Term{classifying-topos}{分类拓扑斯} $B_G$，得到\Term{topos}{拓扑斯} $B_{\mathscr G}^T$，它是\Term{subcategory}{子范畴} $B^T(i)$ 的\Term{inductive-limit}{归纳极限}，而 $B^T(i)$（非\Term{canonical}{典范}地）\Term{equivalence}{等价}于\Term{classifying-topos}{分类拓扑斯} $B_{G_i}$。证明，\Term{topos}{拓扑斯} $B_{\mathscr G}^T$ 有\Term{fiber-functor}{纤维函子}，当且仅当\Term{gerbe}{凝集} $T$“中性”；为此在 $B_{\mathscr G}^T$ 的\Term{fiber-functor}{纤维函子}\Term{category}{范畴}与 $T$ 在 $I$ 上的\Term{section}{截面}\Term{category}{范畴}之间建立\Term{equivalence-of-categories}{范畴等价}。若各 $G_i$ 交换，从而 $I$ 上以 $G$ 为带的\Term{gerbe}{凝集}由 $\varprojlim_I^{(2)}G_i=H^2(I,G)$ 分类（同上），则推出一个形如 $B_{\mathscr G}^T$ 的（并非“空”的（2.2））\Term{topos}{拓扑斯}没有点的例子。
\SourcePage{249}
（取一个 $\varprojlim_I^{(2)}G_i\ne0$ 的例子。）
\end{SGAresult}

\textbf{7.3。\Term{topos}{拓扑斯} $\widehat C$ 的点；其点\Term{category}{范畴}不\Term{equivalence}{等价}于\Term{small-category}{小范畴}的 $U$-\Term{topos}{拓扑斯} $\widehat C$ 的例子。} 设 $E$ 为 $U$-\Term{topos}{拓扑斯}，$(\varphi_i)_{i\in I}$ 为 $E$ 上的\Term{filtered-category}{滤过}\Term{fiber-functor}{纤维函子}族。由\Term{functor}{函子} $\varinjlim$ 的\Term{exactness}{正合性}（第 I 讲），\Term{functor}{函子} $\varphi=\varinjlim\varphi_i$ 仍为\Term{fiber-functor}{纤维函子}：\Term{fiber-functor}{纤维函子}的任意\Term{filtered-category}{滤过}\Term{inductive-limit}{归纳极限}都是\Term{fiber-functor}{纤维函子}。因此，\Term{category}{范畴} $\operatorname{Fib}(E)$ 有 $U$-\Term{filtered-category}{滤过}\Term{inductive-limit}{归纳极限}（且包含\Term{functor}{函子} $\operatorname{Fib}(E)\to\mathscr{H}om(E,(U\text{-}\mathrm{Ens}))$ 与之交换）；换言之，\Term{category}{范畴} $\operatorname{Point}(E)$ 有 $U$-\Term{filtered-category}{滤过}\Term{projective-limit}{投射极限}。此事实在习题 7.1.10 中已用来给出“大”的\Term{fiber-functor}{纤维函子}\Term{category}{范畴}的例子。

利用 (6.8.6.1)，可以用形如 $E=\widehat C$、$C\in U$ 的\Term{topos}{拓扑斯}构造明显更简单的例子。它立即给出 $\operatorname{Fib}(\widehat C)$ 不\Term{equivalence}{等价}于属于 $U$ 的\Term{category}{范畴}的例子，即此 $U$-\Term{category}{范畴}的\Term{object}{对象}\Term{isomorphism}{同构}类之集的\Term{cardinal}{基数}不属于 $U$。这意味着\Term{category}{范畴} $\operatorname{Pro}(C)$ 不\Term{equivalence}{等价}于属于 $U$ 的\Term{category}{范畴}。例如只需取 $D=C^\circ$ 为形如 $[0,n]$ 的\Term{finite}{有限}集的\Term{category}{范畴}，其中 $n\geq0$ 为整数；此时 $\operatorname{Ind}(D)\simeq\operatorname{Pro}(C)^\circ$ \Term{equivalence}{等价}于\Term{nonempty}{非空}\Term{set}{集合}的\Term{category}{范畴}。此时\Term{topos}{拓扑斯} $E$ 就是熟知的余\Term{simplicial-set}{单纯集合}\Term{category}{范畴}。也可以取 $C$ 为\Term{nonempty}{非空}\Term{finite}{有限}集的\Term{category}{范畴}，可见\Term{simplicial-set}{单纯集合}的\Term{topos}{拓扑斯} $\widehat C$ 上的点\Term{equivalence-of-categories}{范畴等价}于\Term{profinite-set}{投有限集}的\Term{category}{范畴}，或者说\Term{equivalence}{等价}于紧致全不\Term{connected}{连通}\Term{space}{空间}的\Term{category}{范畴}。

\textbf{7.4。没有点的\Term{nonempty}{非空}\Term{topos}{拓扑斯}。} 下例归于 P. Deligne。（另一个例子见 7.2.6 d)。）取带有测度 $\mu$ 的\Term{compact-space}{紧空间} $K$，及 $K$ 的可测\Term{subset}{子集}模去零测集所得的\Term{ordered-set}{有序集} $\mathcal U$。把 $\mathcal U$ 看成\Term{category}{范畴} $\mathbf U$，其中 $\operatorname{ob}\mathbf U=\mathcal U$，\Term{morphism}{态射}为 $\mathcal U$ 的元素之间的“包含\Term{morphism}{态射}”。在 $\mathbf U$ 上取\Term{pretopology}{预拓扑}，使 $\operatorname{Cov}(E)$（$E\in\mathcal U$）由所有可数族 $E_i$ 构成，各 $E_i$ 以 $E$ 为上界，且 $E$ 在模去零测集的意义下为这些 $E_i$ 的并。于是使 $\mathbf U$ 成为\Term{site}{景}，\Term{derivation}{导出}\Term{topos}{拓扑斯} $\operatorname{Top}(\mu)=\mathbf U^\sim$，它以\Term{final-object}{终对象}的\Term{subobject}{子对象}之集为\Term{generating-family}{生成族}（这个\Term{topos}{拓扑斯}似乎未曾引起概率学家的注意）。此\Term{topos}{拓扑斯}是“\Term{empty-topos}{空拓扑斯}”（2.2），当且仅当 $\mu=0$。另一方面，这个\Term{topos}{拓扑斯}的点\Term{equivalence-of-categories}{范畴等价}于满足 $\mu(\{x\})\ne0$ 的 $x\in K$ 所构成\Term{set}{集合}定义的\Term{discrete-category}{离散范畴}（证明留给读者）。因此，若 $K$ 没有这样的点，这个\Term{category}{范畴}便为空；
\SourcePage{250}
例如 $K$ 为实直线上的单位闭区间，配备由 Lebesgue 测度诱导的测度时便如此。

\begin{SGAresult}{习题}{7.5}
（\Term{karoubian-category}{Karoubi 范畴}与\Term{topos-morphism}{拓扑斯态射}。）\footnote{比较第 I 讲 8.7.8。}

 a) 设 $C$ 为\Term{category}{范畴}，$X$ 为 $C$ 的\Term{object}{对象}，$p$ 为 $X$ 的\Term{endomorphism}{自同态}。若 $p^2=p$，则称 $p$ 为\Term{projector}{投影子}。证明，若 $p$ 为\Term{projector}{投影子}，则 $\operatorname{Ker}(\operatorname{id}_X,p)$ \Term{representable}{可表}当且仅当 $\operatorname{Coker}(\operatorname{id}_X,p)$ \Term{representable}{可表}，且这样得到的 $C$ 的两个\Term{object}{对象}\Term{canonical}{典范}\Term{isomorphism}{同构}。此时说\Term{projector}{投影子} $p$ 有像，并称 $\operatorname{Ker}(\operatorname{id}_X,p)\simeq\operatorname{Coker}(\operatorname{id}_X,p)$ 为\Term{projector}{投影子} $p$ 的像；依上下文将它等同于 $C$ 的一个\Term{subobject}{子对象}或\Term{quotient-object}{商对象}。若某\Term{object}{对象}\Term{isomorphism}{同构}于 $X$ 中某个适当\Term{projector}{投影子} $p$ 的 $\operatorname{Im}p$，则称它为 $C$ 的\Term{object}{对象} $X$ 的一个\Term{direct-factor}{直因子}。若 $C$ 的\Term{object}{对象}中的每个\Term{projector}{投影子}都有像，则称 $C$ 为有\Term{direct-factor}{直因子}的\Term{category}{范畴}，或\Term{karoubian-category}{Karoubi 范畴}。若\Term{functor}{函子} $F:C\to C'$ 与\Term{kernel}{核}或\Term{cokernel}{余核}交换，则 $F$ 把有像的\Term{projector}{投影子}变为有像的\Term{projector}{投影子}。

 b) 证明，对于任意\Term{category}{范畴} $C$，可以找到从 $C$ 到一个 \Term{karoubian-category}{Karoubi 范畴}（在\Term{equivalence}{等价}意义下确定）的\Term{functor}{函子} $\varphi:C\to\operatorname{kar}(C)$，使对任意 \Term{karoubian-category}{Karoubi 范畴} $C'$，\Term{functor}{函子}
\[f\longmapsto f\circ\varphi:\mathscr{H}om(\operatorname{kar}(C),C')\longrightarrow\mathscr{H}om(C,C')\]
为\Term{equivalence-of-categories}{范畴等价}；称 $\operatorname{kar}(C)$ 为\Term{category}{范畴} $C$ 的 \Term{karoubi-envelope}{Karoubi 包络}。（提示：取 $\operatorname{ob}\operatorname{kar}(C)$ 为二元组 $(X,p)$ 之集，其中 $X\in\operatorname{ob}C$，$p$ 为 $X$ 中的\Term{projector}{投影子}；取 $\operatorname{Hom}((X,p),(Y,q))$ 为 $\operatorname{Hom}(X,Y)$ 中满足 $f=qfp$ 的 $f:X\to Y$ 所组成的\Term{subset}{子集}。）证明 $\varphi$ \Term{fully-faithful-functor}{全忠实}。

 c) 设 $F:E\to E'$ 为从一个\Term{category}{范畴}到另一个\Term{category}{范畴}的\Term{functor}{函子}。证明，若 $F$ 与某类型的\Term{inductive-limit}{归纳极限}或\Term{projective-limit}{投射极限}交换，则 $F$ 的任意\Term{direct-factor}{直因子}也如此。特别地，若 $E,E'$ 为 $U$-\Term{topos}{拓扑斯}，且 $F$ 是某个\Term{topos-morphism}{拓扑斯态射} $E'\to E$ 的\Term{inverse-image-functor}{逆像函子}，则 $F$ 的任意\Term{direct-factor}{直因子}也如此；因而\Term{category}{范畴} $\mathscr{H}om\mathrm{top}(E',E)$ 为 \Term{karoubian-category}{Karoubi 范畴}，特别地，\Term{category}{范畴} $\operatorname{Point}(E)$ 为 \Term{karoubian-category}{Karoubi 范畴}。

 d) 证明，对于任意 $U$-\Term{category}{范畴} $C$，\Term{category}{范畴} $\operatorname{Ind}(C)$ 为 \Term{karoubian-category}{Karoubi 范畴}（其中 $\operatorname{Ind}(C)$ 由以\Term{filtered-category}{滤过}\Term{preordered-set}{预序集} $I\in U$ 为指标的 $C$ 中\Term{inductive-system}{归纳系}构成）。（若 $C\in U$，例如可利用事实（7.3）：$\operatorname{Ind}(C)$ \Term{equivalence}{等价}于\Term{category}{范畴} $\operatorname{Fib}(\widehat{C^\circ})=\operatorname{Point}(\widehat{C^\circ})^\circ$。）由此给出 $\operatorname{kar}(C)$ 的另一构造：它为 $\operatorname{Ind}(C)$ 中由 $C$ 的\Term{object}{对象} $X$ 的\Term{projector}{投影子}在 $\operatorname{Ind}(C)$ 中的像构成的\Term{full-subcategory}{满子范畴}。
\end{SGAresult}

\begin{SGAresult}{习题}{7.6}
（本质\Term{topos-morphism}{拓扑斯态射}，\Term{essential-point}{本质点}。）
\SourcePage{251}
 a) 设 $f:E\to F$ 为\Term{topos-morphism}{拓扑斯态射}。证明以下条件\Term{equivalence}{等价}：

 i) $f_!$ 存在，即 $f^*$ 有\Term{left-adjoint}{左伴随}。

 ii) $f^*$ 与 $U$-\Term{projective-limit}{投射极限}交换。

 ii') $f^*$ 与 $U$-积交换。

此时称 $f$ 为从\Term{topos}{拓扑斯} $E$ 到\Term{topos}{拓扑斯} $E'$ 的\Term{essential-morphism}{本质态射}。证明，若 $f$ 满足这些条件，则 $f$ 的任意\Term{direct-factor}{直因子}也如此。（利用 1.8 和 7.5 c)。）

 b) 若 $E$ 的点 $p:P\to E$ 满足 $p_!$ 存在，则称其为\Term{topos}{拓扑斯} $E$ 的\Term{essential-point}{本质点}。证明，若 $E=\operatorname{Top}(X)$，其中 $X$ 为\Term{topological-space}{拓扑空间}，且 $p$ 为由 $x\in X$ 定义的 $E$ 的点，则 $p$ 本质当且仅当 $x$ 有最小开\Term{neighborhood}{邻域} $U$（若 $X$ 的点都是闭的，这就是说 $x$ 为 $X$ 的\Term{isolated-point}{孤立点}）。\Term{topos-morphism}{拓扑斯态射} $f:E\to F$ 为\Term{essential-morphism}{本质态射}（a)）的一个必要条件是，$f$ 把\Term{essential-point}{本质点}变为\Term{essential-point}{本质点}；当 $E$ 有足够多的\Term{essential-point}{本质点}时（即这些点所成的族\Term{conservative-functor}{保守}），这个条件也充分（见 d)）。

 c) \Term{topos}{拓扑斯} $E$ 的点 $p$ 本质的充要条件是，\Term{correspondence}{对应}\Term{fiber-functor}{纤维函子}可由 $E$ 的\Term{object}{对象} $X$ \Term{represented}{表出}。给定 $E$ 的\Term{object}{对象} $X$，它所\Term{represented}{表出}的\Term{covariant-functor}{协变函子} $\varphi:E\to(\mathrm{Ens})$ 为\Term{fiber-functor}{纤维函子}（即定义 $E$ 的一个点，此点必本质）的充要条件是：$X$ \Term{connected-nonempty}{连通非空}（4.3.5）且\Term{projective-module}{投射}（即\Term{functor}{函子} $\varphi$ 把\Term{epimorphism}{满态射}变为\Term{epimorphism}{满态射}）。（提示：先证明 $\varphi$ 与和交换的充要条件是 $X$ \Term{connected-nonempty}{连通非空}，再利用判据 4.9.4。）由此给出\Term{category}{范畴} $\operatorname{Pointess}(E)$ 与 $E$ 中由\Term{connected-nonempty}{连通非空}\Term{projective-object}{投射对象}组成的\Term{full-subcategory}{满子范畴} $P$ 之间的\Term{equivalence}{等价}。

 d) 证明，由 $E$ 的\Term{topology}{拓扑}在 $P$ 上诱导的\Term{topology}{拓扑}为\Term{chaotic-topology}{混沌拓扑}。由此推出关于 $E$ 的以下条件\Term{equivalence}{等价}（归于 J. E. Roos [12 c)，命题 1]）：(i) $E$ 的\Term{essential-point}{本质点}族\Term{conservative-functor}{保守}；(ii) $E$ 中由\Term{connected-nonempty}{连通非空}\Term{projective-object}{投射对象}构成的\Term{full-subcategory}{满子范畴} $P$ 是生成的；(iii) $E$ \Term{equivalence}{等价}于形如 $\widehat C$ 的\Term{topos}{拓扑斯}，其中 $C$ 为\Term{equivalence}{等价}于属于 $U$ 的\Term{category}{范畴}的\Term{category}{范畴}（或者也可取 $C\in U$）。（利用 c)，以及若 $P$ 生成，则自然\Term{functor}{函子} $E\to P^\sim$ 为\Term{equivalence-of-categories}{范畴等价}这一事实。）

 e) \Term{topos}{拓扑斯} $E$ 的\Term{essential-point}{本质点}\Term{category}{范畴}为 \Term{karoubian-category}{Karoubi 范畴}（见 7.5 c)）。设 $C$ 为\Term{equivalence}{等价}于属于 $U$ 的\Term{category}{范畴}的\Term{category}{范畴}。证明，\Term{canonical}{典范}的\Term{fully-faithful-functor}{全忠实}\Term{functor}{函子}（4.6.2）$C\to\operatorname{Point}(\widehat C)$ 经\Term{functor}{函子}
\[C\longrightarrow\operatorname{Pointess}(\widehat C)\]
\Term{resolution}{分解}，该\Term{functor}{函子}使 $\operatorname{Pointess}(\widehat C)$ 成为 $C$ 的 \Term{karoubi-envelope}{Karoubi 包络}（7.5 b)）。特别地，此\Term{functor}{函子}为\Term{equivalence-of-categories}{范畴等价}，当且仅当\Term{category}{范畴} $C$ 为
\SourcePage{252}
\Term{karoubian-category}{Karoubi 范畴}。（提示：利用 c)，证明 $\widehat C$ 的每个\Term{isolated-point}{孤立点}都\Term{isomorphism}{同构}于某个 $X\in\operatorname{ob}C$ 的\Term{direct-factor}{直因子}。）

 f) 设 $C,C'$ 为\Term{equivalence}{等价}于属于 $U$ 的\Term{category}{范畴}的两个\Term{category}{范畴}。证明，\Term{canonical}{典范}的\Term{fully-faithful-functor}{全忠实}\Term{functor}{函子}（4.6.2）
\[\mathscr{H}om(C,C')\longrightarrow\mathscr{H}om\mathrm{top}(\widehat C,\widehat{C'})\]
取值于由本质\Term{topos-morphism}{拓扑斯态射}（a)）组成的\Term{full-subcategory}{满子范畴} $\mathscr{H}om\mathrm{topess}(\widehat C,\widehat{C'})$，且诱导\Term{functor}{函子}
\[\mathscr{H}om(C,C')\longrightarrow\mathscr{H}om\mathrm{topess}(\widehat C,\widehat{C'})\]
为\Term{equivalence-of-categories}{范畴等价}，当且仅当 $C$ 为空或 $C'$ 为 \Term{karoubian-category}{Karoubi 范畴}。（利用下述 g)。）

 g) 设 $C$ 为\Term{equivalence}{等价}于属于 $U$ 的\Term{category}{范畴}的\Term{category}{范畴}，$E$ 为\Term{topos}{拓扑斯}。在本质\Term{topos-morphism}{拓扑斯态射} $\widehat C\to E$ 的\Term{category}{范畴} $\mathscr{H}om\mathrm{topess}(\widehat C,E)$ 与\Term{category}{范畴} $\mathscr{H}om(C,\operatorname{Pointess}(E))$ 之间定义一个\Term{equivalence}{等价}。（利用 b) 及 $\widehat C$ 有足够多的\Term{essential-point}{本质点}这一事实。）

 h) 设 $C$ 为\Term{finite}{有限}\Term{category}{范畴}。证明，$\widehat C$ 的每个点均为\Term{essential-point}{本质点}，且 $\operatorname{Point}(\widehat C)$ \Term{equivalence}{等价}于\Term{finite}{有限}\Term{category}{范畴}。（注意 $C$ 的 \Term{karoubi-envelope}{Karoubi 包络}\Term{finite}{有限}，并利用 (6.8.6.1)，将问题归结为证明：若 $C$ 为\Term{finite}{有限} \Term{karoubian-category}{Karoubi 范畴}，则\Term{canonical}{典范}\Term{functor}{函子} $C\to\operatorname{Pro}(C)$ 为\Term{equivalence-of-categories}{范畴等价}。）利用 b) 推出，每个\Term{topos-morphism}{拓扑斯态射} $\widehat{C'}\to\widehat C$ 都本质，且 $\mathscr{H}om(C,C')\to\operatorname{Hom}(\widehat C,\widehat{C'})$ 为\Term{equivalence}{等价}，当且仅当 $C$ 为空或 $C'$ 为 \Term{karoubian-category}{Karoubi 范畴}。

 i) 设 $X$ 为\Term{topological-space}{拓扑空间}，$f:\operatorname{Top}(X)\to P$ 为由 $X$ 到单点\Term{space}{空间}的\Term{continuous-map}{连续映射}\Term{derivation}{导出}的\Term{topos-morphism}{拓扑斯态射}。证明，$f$ 本质当且仅当\Term{space}{空间} $X$ \Term{locally-connected}{局部连通}，即满足下述条件：对每个 $x\in X$ 及 $x$ 的任意开\Term{neighborhood}{邻域} $U$，$x$ 在 $U$ 中的\Term{connected-component}{连通分支}为 $x$ 的\Term{neighborhood}{邻域}；或者说，对 $X$ 的每个\Term{open-subset}{开集} $U$，$U$ 的\Term{connected-component}{连通分支}都是开的。（推广见 8.7 b)。）
\end{SGAresult}

\begin{SGAresult}{习题}{7.7}
（\Term{classifying-topos}{分类拓扑斯}的不寻常的点。）\footnote{另见 7.2.6 d)。}

 a) 设 $G$ 为\Term{monoid}{幺半群}。$G$ 含有满足 $g\ne e$、$g^2=g$ 的元素 $g$，当且仅当\Term{classifying-topos}{分类拓扑斯} $B_G$ 有一个不\Term{isomorphism}{同构}于\Term{trivial-bundle}{平凡}点（\Term{correspondence}{对应}于“遗忘 $G$ 的\Term{action}{作用}”的\Term{fiber-functor}{纤维函子}）的\Term{essential-point}{本质点}。（利用 7.6 e)。）

 b) 设 $G$ 为非负整数的加法\Term{monoid}{幺半群}。证明，\Term{classifying-topos}{分类拓扑斯} $B_G$ 的每个\Term{essential-point}{本质点}都\Term{isomorphism}{同构}于\Term{trivial-bundle}{平凡}点。构造 $B_G$ 的一个不\Term{isomorphism}{同构}于\Term{trivial-bundle}{平凡}点的点。证明，\Term{category}{范畴} $\operatorname{Point}(B_G)$ 不\Term{equivalence}{等价}于属于 $U$ 的\Term{category}{范畴}。
\end{SGAresult}

\SourcePage{253}
\begin{SGAresult}{习题}{7.8}
（$\operatorname{Point}(E)$ 上的\Term{topology}{拓扑}，以及与\Term{ordered-set}{有序集}\Term{correspondence}{对应}的\Term{topos}{拓扑斯}。）

 a) 设 $E$ 为\Term{topos}{拓扑斯}。以 $\operatorname{Point}'E$ 表示 $E$ 的点在\Term{isomorphism}{同构}意义下的类所组成的\Term{set}{集合}。对 $E$ 的每个\Term{open-subset}{开集} $U$，令 $U'$ 为满足 $U_p\ne\varnothing$ 的 $E$ 的点 $p$ 的类之集。证明，$U\mapsto U'$ 是从 $E$ 的\Term{open-subset}{开集}的\Term{ordered-set}{有序集} $\operatorname{Ouv}(E)$ 到 $X=\operatorname{Point}(E)$ 的幂集的递增映射，且此映射与\Term{finite-infimum}{有限下确界}及\Term{arbitrary-supremum}{任意上确界}交换。由此推出，形如 $U'$ 的 $X$ 的\Term{subset}{子集}之集定义 $X$ 上的一个\Term{topology}{拓扑}（称为\Term{topos}{拓扑斯} $E$ 的点的类之集上的\Term{canonical-topology}{典范拓扑}）。证明，若 $E$ 有\Term{enough-points}{足够多的点}，则映射 $U\mapsto U'$ 是\Term{ordered-set}{有序集} $\operatorname{Ouv}(E)$ 到\Term{ordered-set}{有序集} $\operatorname{Ouv}(X)$ 的\Term{isomorphism}{同构}。

 b) 设 $\mathcal O\in U$ 为有\Term{finite-infimum}{有限下确界}及\Term{arbitrary-supremum}{任意上确界}的\Term{ordered-set}{有序集}，从而定义具有\Term{finite-product}{有限积}及任意和的\Term{category}{范畴} $\operatorname{cat}(\mathcal O)$。若 $U$ 为各 $U_i$ 的\Term{supremum}{上确界}，则称具有同一目标 $U$ 的\Term{morphism}{态射}族 $U_i\to U$ 为\Term{covering-family}{覆盖族}。假定 $\operatorname{cat}(\mathcal O)$ 中的和是\Term{universal}{普遍}的，即 $\mathcal O$ 中\Term{arbitrary-supremum}{任意上确界}与\Term{infimum}{下确界}交换，便由此在 $\operatorname{cat}(\mathcal O)$ 上定义一个使它成为\Term{site}{景}的\Term{topology}{拓扑}，从而得到\Term{topos}{拓扑斯} $\operatorname{cat}(\mathcal O)^\sim$，也简记为 $\mathcal O^\sim$。证明，$I$ \Term{canonical}{典范}\Term{isomorphism}{同构}于 $\operatorname{Ouv}(\mathcal O^\sim)$。

证明，\Term{category}{范畴} $\mathscr{H}om\mathrm{top}(E,\mathcal O^\sim)$ \Term{equivalence}{等价}于如下\Term{ordered-set}{有序集}所\Term{correspondence}{对应}的\Term{category}{范畴}：从 $\mathcal O$ 到 $\operatorname{Ouv}(E)$ 且与\Term{finite-infimum}{有限下确界}及\Term{arbitrary-supremum}{任意上确界}交换的\Term{ordered-set}{有序集}\Term{morphism}{态射}所成的\Term{ordered-set}{有序集}。（利用判据 4.9.4。）若 $\mathcal O'$ 为满足与 $\mathcal O$ 相同条件的\Term{ordered-set}{有序集}，推出 $\mathscr{H}om\mathrm{top}(\mathcal O^\sim,\mathcal O'^\sim)$ \Term{equivalence}{等价}于由所有与\Term{finite-infimum}{有限下确界}及\Term{arbitrary-supremum}{任意上确界}交换的\Term{ordered-set}{有序集}\Term{morphism}{态射} $\mathcal O\to\mathcal O'$ 构成的\Term{ordered-set}{有序集}所\Term{correspondence}{对应}的\Term{category}{范畴}。

 c) 证明，\Term{topos}{拓扑斯} $E$ \Term{equivalence}{等价}于形如 $\mathcal O^\sim$ 的\Term{topos}{拓扑斯}（$\mathcal O$ 如 b) 中所述）的充要条件为：$E$ 的\Term{open-subset}{开集}族在 $E$ 中生成。$E$ \Term{equivalence}{等价}于某个 $\operatorname{Top}(X)$（$X\in U$）的充要条件为：它满足上述条件且有\Term{enough-points}{足够多的点}。假定第一个条件成立，利用 b) 的最后一个断言，用\Term{ordered-set}{有序集} $\mathcal O=\operatorname{Ouv}(E)$ 的结构表达第二个条件。

 d) 定义\Term{canonical-topology}{典范拓扑}斯\Term{morphism}{态射}
\[E\longrightarrow\operatorname{Ouv}(E)^\sim,\]
使它对于从 $E$ 到由其\Term{open-subset}{开集}生成的\Term{topos}{拓扑斯}（见 c)）的\Term{morphism}{态射}具有\Term{universal}{普遍}性（在\Term{equivalence}{等价}意义下）。

 e) 设 $X'$ 为 $X=\operatorname{Point}(E)$ 的一个小\Term{subset}{子集}，配备 $X$ 的\Term{topology}{拓扑}所诱导的\Term{topology}{拓扑}。定义\Term{canonical-topology}{典范拓扑}斯\Term{morphism}{态射}
\[\operatorname{Ouv}(E)^\sim\longrightarrow\operatorname{Top}(X'),\]
\SourcePage{254}
使它当 $E$ 的点族 $X'$ \Term{conservative-functor}{保守}时为\Term{equivalence}{等价}。由此得到\Term{canonical-topology}{典范拓扑}斯\Term{morphism}{态射}
\[E\longrightarrow\operatorname{Top}(X').\]
对于从\Term{topos}{拓扑斯} $E$ 到形如 $\operatorname{Top}(Y)$ 的\Term{topos}{拓扑斯}的\Term{morphism}{态射}，给出此\Term{topos-morphism}{拓扑斯态射}的\Term{universal}{普遍}刻画（在\Term{equivalence}{等价}意义下）。（在这里注意，到\Term{equivalence}{等价}为止，$\operatorname{Top}(X')$ 与选取的小\Term{conservative-functor}{保守}\Term{fiber-functor}{纤维函子}族无关；或\Term{equivalence}{等价}地（4.2），$X'_{\mathrm{sob}}$ 的\Term{homeomorphism}{同胚}类型与这样的族的选择无关。）

 f) 假定\Term{category}{范畴} $\operatorname{Point}(E)$ 不\Term{equivalence}{等价}于\Term{small-category}{小范畴}（见 7.3），且 $E$ 有一个小的\Term{conservative-functor}{保守}\Term{fiber-functor}{纤维函子}族。证明，若把这样的族取得\Term{large-relative-to-cardinal}{足够大}，则它所定义的 $\operatorname{Point}(E)$ 的\Term{subset}{子集} $X'$ 对于\Term{induced-topology}{诱导拓扑}不是\Term{sober-space}{清醒拓扑空间}。

 g) 关于与习题 7.6 的关系，见 8。
\end{SGAresult}

\SGAsection{8}{局部化。拓扑斯的开集}
\textbf{8.1。} 设 $C$ 为\Term{category}{范畴}，$T(X)$ 为涉及 $C$ 的\Term{object}{对象} $X$ 的关系。若对 $C$ 的任意\Term{morphism}{态射} $X'\to X$，关系 $T(X)$ 蕴含 $T(X')$，则称关系 $T(X)$（关于 $X$）在\Term{base-change}{换基}下\Term{stable}{稳定}。这\Term{equivalence}{等价}于说，$C$ 中由满足 $T(X)$ 的\Term{object}{对象} $X$ 所成的\Term{full-subcategory}{满子范畴} $R$ 为筛（第 I 讲 4.1）。

\textbf{注记 8.1.1。} 假定 $C$ 是 $U$-\Term{category}{范畴}（第 I 讲 1.1），从而 $C$ 可等同于 $\widehat C$ 的一个\Term{full-subcategory}{满子范畴}（第 I 讲 1.4）。有时，把 $\operatorname{ob}C$ 上的关系 $T$ 扩充为涉及 $\widehat C$ 的\Term{object}{对象} $F$ 的关系 $\widehat T$ 很方便：以 $\widehat T(F)$ 表示下述关系：对 $\widehat C$ 中任意\Term{arrow}{箭头} $X\to F$，其中 $X\in\operatorname{ob}C$，都有 $T(X)$。关系 $\widehat T(F)$ 显然仍关于 $F$ 在\Term{base-change}{换基}下\Term{stable}{稳定}。仍以 $R$ 表示由 $C$ 的筛 $R$ 所定义的 $\widehat C$ 的\Term{final-object}{终对象} $e$ 的\Term{subobject}{子对象}（第 I 讲 4.2.1），则关系 $\widehat T(F)$ 也\Term{representable}{可表}达为 $\operatorname{Hom}_{\widehat C}(F,R)\ne\varnothing$，即唯一的\Term{morphism}{态射} $F\to e$ 经 $e$ 的\Term{subobject}{子对象} $R$ \Term{resolution}{分解}。

\textbf{8.2。} 现在假定 $C$ 为\Term{site}{景}。若对于 $C$ 的\Term{object}{对象} $X$ 的任意\Term{covering-family}{覆盖族} $(f_i:X_i\to X)_{i\in I}$，关系“对每个 $i\in I$ 均有 $T(X_i)$”蕴含 $T(X)$，则称关于自变量 $X\in\operatorname{ob}C$ 的关系 $T(X)$ 在\Term{descent}{下降}下\Term{stable}{稳定}。若 $T$ 在\Term{base-change}{换基}下\Term{stable}{稳定}（8.1）且在\Term{descent}{下降}下\Term{stable}{稳定}，则称 $T$ 为具有\Term{local-nature}{局部性质}的关系。当 $C$ 为 $U$-\Term{category}{范畴}，并且已经假定 $T$ 在\Term{base-change}{换基}下\Term{stable}{稳定}，即它可由 $C$ 的一个筛 $R$ 定义，且把 $R$ 等同于 $C$ 上\Term{final-presheaf}{终预层}的\Term{subpresheaf}{子预层}时，立即可见 $C$ 具有\Term{local-nature}{局部性质}，当且仅当 $R$ 是\Term{sheaf}{层}。因此，关于一个自变量 $X\in\operatorname{ob}C$ 的具有\Term{local-nature}{局部性质}的关系，恰与 $C$ 的\Term{final-sheaf}{终层}的\Term{subsheaf}{子层}，即\Term{topos}{拓扑斯} $C^\sim$ 的\Term{final-object}{终对象}的\Term{subobject}{子对象}\Term{correspondence}{对应}。
\SourcePage{255}
所以它们实质上只依赖于\Term{site}{景} $C$ 所定义的\Term{topos}{拓扑斯} $C^\sim$（如同与\Term{site}{景}\Term{correspondence}{对应}的所有重要概念一样！）。用\Term{final-sheaf}{终层}的\Term{subsheaf}{子层} $R$ 表述时，关系 $T(X)$ 写成 $\operatorname{Hom}(\varepsilon(X),R)\ne\varnothing$。它\Term{canonical-extension}{典范延拓}为\Term{topos}{拓扑斯} $C^\sim$ 上的关系 $T^\sim(F):\operatorname{Hom}(F,R)\ne\varnothing$。

\textbf{注记 8.2.1。} 沿用 8.1.1 引入的记号，对 $C$ 上的\Term{presheaf}{预层} $F$，因为 $\operatorname{Hom}(F,R)\simeq\operatorname{Hom}(\underline F,R)$，关系 $\widehat T(F)$ \Term{equivalence}{等价}于关系 $\widehat T(\underline F)$。

\begin{SGAresult}{定义}{8.3}
$U$-\Term{topos}{拓扑斯} $E$ 的\Term{final-object}{终对象} $e_E$ 的任意\Term{subobject}{子对象}称为 $E$ 的\Term{open-topos}{开集}；若 $X$ 为 $E$ 的\Term{object}{对象}，有时把\Term{induced-topos}{诱导拓扑斯} $E_{/X}$ 的\Term{open-subset}{开集}，即 $X$ 的任意\Term{subobject}{子对象}，称为 $X$ 的\Term{open-subset}{开集}。

$U$-\Term{site}{景} $C$ 的\Term{open-subset}{开集}，是指\Term{correspondence}{对应} $U$-\Term{topos}{拓扑斯} $C^\sim$ 的\Term{open-subset}{开集}。
\end{SGAresult}

注意，由于 $E$ 的\Term{final-object}{终对象}彼此\Term{canonical}{典范}\Term{isomorphism}{同构}，8.3 中由 $e_E$ 的选择引入的歧义无关紧要；也可更内在地把 $E$ 的\Term{open-subset}{开集}定义为这样的\Term{full-subcategory}{满子范畴} $R$：对 $E$ 的\Term{object}{对象} $F$，关系 $F\in\operatorname{ob}R$ 具有\Term{local-nature}{局部性质}，不过实际使用时这不太方便。同样，$U$-\Term{site}{景} $C$ 的\Term{open-subset}{开集}与 $C$ 的那些使关于 $C$ 的\Term{object}{对象}的关系 $F\in\operatorname{ob}R$ 具有\Term{local-nature}{局部性质}的\Term{full-subcategory}{满子范畴}一一\Term{correspondence}{对应}；所以此概念实质上与\Term{universe}{宇宙} $U$ 的选择无关。最后，由 8.2 中所述，\Term{topos}{拓扑斯}（或\Term{site}{景}）上具有\Term{local-nature}{局部性质}的关系，实质上与该\Term{topos}{拓扑斯}（或\Term{site}{景}）的\Term{open-subset}{开集}是同一回事。

\textbf{8.4。\Term{open-subset}{开集}的例子。}

\textbf{8.4.1。} 设 $X$ 为\Term{topological-space}{拓扑空间}，把 $\operatorname{Top}(X)$（2.1）解释为 $X$ 上\Term{etale-space}{平展空间}的\Term{category}{范畴}。因为 $\operatorname{Top}(X)$ 中的\Term{monomorphism}{单态射}显然为\Term{injection}{单射}映射，故\Term{topos}{拓扑斯} $T(X)$ 的\Term{open-subset}{开集}等同于\Term{space}{空间} $X$ 的\Term{open-subset}{开集}。更一般地，\Term{topos}{拓扑斯} $\operatorname{Top}(X,G)$（2.4）的\Term{open-subset}{开集}等同于 $X$ 的在 $G$ \Term{action}{作用}下不变的\Term{open-subset}{开集}。

\textbf{8.4.2。} $U$-\Term{topos}{拓扑斯} $E$ 的\Term{open-subset}{开集}构成一个 $U$-\Term{small-set}{小集合}（第 I 讲 7.4），按包含排序，且有\Term{arbitrary-supremum}{任意上确界}和\Term{finite-infimum}{有限下确界}。它以\Term{initial-object}{始对象}（或“空\Term{object}{对象}”）$\varnothing_E$ 为最小元素，以 $E$ 的\Term{final-object}{终对象} $e_E$ 为最大元素。只有当 $E$ 为“\Term{empty-topos}{空拓扑斯}”（2.2）时，这两个元素才相同。

\textbf{8.4.3。} 设 $G$ 为 $E$ 中的\Term{monoid}{幺半群}，从而有\Term{topos}{拓扑斯} $B_G$（2.4），即 $E$ 中带有 $G$ 的左\Term{action}{作用}的\Term{object}{对象}的\Term{category}{范畴}。$B_G$ 的\Term{final-object}{终对象} $e_G$ 为带有 $G$ 的\Term{trivial-action}{平凡作用}的 $E$ 的\Term{final-object}{终对象} $e_E$。所以 $B_G$ 的\Term{open-subset}{开集}\Term{ordered-set}{有序集}\Term{canonical}{典范}\Term{isomorphism}{同构}于 $E$ 的\Term{open-subset}{开集}\Term{category}{范畴}。特别地，若 $E$ 为\Term{punctual-topos}{点拓扑斯}，因而 $G$ 为普通\Term{monoid}{幺半群}，则 $B_G$ 的\Term{open-subset}{开集}之集恰有两个元素，即 $\varnothing_{B_G}$ 和 $e_G$。

\SourcePage{256}
\textbf{8.4.4。} 若 $E$ 为形如 $\widehat C$ 的\Term{topos}{拓扑斯}，其中 $C$ 为 $U$-\Term{category}{范畴}，则由 8.1，$E$ 的\Term{open-subset}{开集}\Term{ordered-set}{有序集}\Term{isomorphism}{同构}于 $C$ 的筛的\Term{ordered-set}{有序集}。

\textbf{8.4.5。} 设 $X$ 为非\Term{discrete}{离散}\Term{sober-space}{清醒拓扑空间}。证明，$\operatorname{TOP}(X)$（2.5）有一些\Term{open-subset}{开集}，不能由 $X$ 的\Term{open-subset}{开集}，即 $\operatorname{Top}(X)$ 的\Term{open-subset}{开集}，通过\Term{canonical}{典范}\Term{morphism}{态射} $\operatorname{TOP}(X)\to\operatorname{Top}(X)$（4.10.1）取\Term{inverse-image}{逆像}得到。

\textbf{8.5。具有\Term{local-nature}{局部性质}的关系的例子。} 对\Term{topos}{拓扑斯} $E$ 的每个\Term{object}{对象} $X$，以 $F\mapsto F_X$ 表示\Term{localization-functor}{局部化函子} $j_X^*:E\to E_{/X}$、$Y\mapsto Y\times X$（5.2）。

\textbf{8.5.1。} 设 $u:F\to G$ 为 $E$ 的\Term{morphism}{态射}。关于自变量 $X\in\operatorname{ob}E$ 的关系
\[\text{$u_X:F_X\to G_X$ 是单态射（或满态射，或同构）}\]
具有\Term{local-nature}{局部性质}。特别地，若 $G$ 为 $E$ 中的\Term{group}{群}，则关系“$G_X$ 为 $X$ 上的\Term{unit-group}{单位群}”具有\Term{local-nature}{局部性质}；\Term{correspondence}{对应}的 $E$ 的\Term{open-subset}{开集}称为\Term{group}{群} $G$ 的\Term{cosupport}{余支集}。

\textbf{8.5.2。} 设 $u,v:F\rightrightarrows G$ 为 $E$ 的两个\Term{morphism}{态射}。关于自变量 $X\in\operatorname{ob}E$ 的关系 $u_X=v_X:F_X\to G_X$ 具有\Term{local-nature}{局部性质}。特别地，若 $G$ 为 $E$ 中的\Term{group}{群}，$u$ 为 $G$ 的\Term{section}{截面}（4.3.6），则关于 $X$ 的关系
\[\text{$E_{/X}$ 中的群 $G_X$ 的截面 $u_X$ 为零}\]
具有\Term{local-nature}{局部性质}；\Term{correspondence}{对应}的 $E$ 的\Term{open-subset}{开集}称为\Term{group}{群} $G$ 的\Term{section}{截面} $u$ 的\Term{cosupport}{余支集}。

\textbf{注记 8.5.3。} 当 $E=\operatorname{Top}(X)$ 时，\Term{group-sheaf}{群层} $G$ 的\Term{cosupport}{余支集}，或这样的\Term{sheaf}{层} $G$ 的\Term{section}{截面}的\Term{cosupport}{余支集}，正是经典术语 [TF] 中 $G$ 的\Term{support}{支集}或 $u$ 的\Term{support}{支集}在 $X$ 中的补\Term{open-subset}{开集}。

\textbf{8.5.4。} 设 $P(f)$ 为关于自变量 $f\in\operatorname{fl}C$ 的关系，$C$ 为\Term{site}{景}。若 $C$ 中\Term{fiber-product}{纤维积}\Term{representable}{可表}，则称关系 $P(f)$ 在目标上（或在基上，或“在下方”）\Term{base-change}{换基}\Term{stable}{稳定}（或\Term{descent}{下降}\Term{stable}{稳定}，或具有\Term{local-nature}{局部性质}），如果对 $C$ 的每个\Term{morphism}{态射} $f:X\to Y$，关于自变量 $Y'\in\operatorname{ob}C_{/Y}$ 的关系
\[\text{“$f$ 沿结构态射 $Y'\to Y$ 换基所得的态射 $f':X'=X\times_Y Y'\to Y'$ 满足 $P(f')$”}\]
在\Term{base-change}{换基}下\Term{stable}{稳定}（或在\Term{descent}{下降}下\Term{stable}{稳定}，或具有\Term{local-nature}{局部性质}）。当 $E$ 的\Term{topology}{拓扑}由一个\Term{pretopology}{预拓扑}定义时，若对 $C$ 的任意\Term{morphism}{态射} $f:X\to Y$ 及 $\operatorname{Cov}(X)$ 的任意族 $(g_i:X_i\to X)_{i\in I}$，关系 $P(f)$ \Term{equivalence}{等价}于关系“对每个 $i\in I$，$P(fg_i)$”，则称关系 $P(f)$ 在源上（或“在上方”）具有\Term{local-nature}{局部性质}。当 $\operatorname{Cov}(X)$ 由
\SourcePage{257}
$C$ 的所有\Term{covering-family}{覆盖族}组成时，此条件也意味着，诱导\Term{site}{景} $C_{/X}$ 的\Term{object}{对象} $X'$ 上的关系
\[P(fg),\quad\text{其中 $g:X'\to X$ 为结构态射}\]
具有\Term{local-nature}{局部性质}。注意，若关系 $P(f)$ 在上方具有\Term{local-nature}{局部性质}，则它当然在下方具有\Term{local-nature}{局部性质}。

\textbf{8.5.5。} 例如取 $C=(\mathrm{Sch})$，即属于 $U$ 的\Term{scheme}{概形}的\Term{category}{范畴}，配备\Term{faithfully-flat}{忠实平坦}\Term{quasi-compact}{拟紧}\Term{topology}{拓扑}（SGA 3 第 IV 讲 6.3）或较粗的\Term{topology}{拓扑}。则\Term{morphism}{态射}的以下每个性质在下方都具有\Term{local-nature}{局部性质}：

\Term{surjection}{满射}、\Term{radicial}{根性}、\Term{universally-open}{普遍开}、\Term{universal}{普遍}\Term{finite}{有限}、\Term{proper}{固有}、\Term{quasi-compact}{拟紧}、\Term{quasi-compact}{拟紧}且\Term{dominant-morphism}{支配}、\Term{universal-homeomorphism}{普遍同胚}、\Term{separated}{分离}、\Term{quasi-separated}{拟分离}、\Term{locally-finite-type}{局部有限型}、\Term{locally-finite-presentation}{局部有限表现}、\Term{finite-type}{有限型}、\Term{finite-presentation}{有限表现}、为\Term{isomorphism}{同构}、为\Term{monomorphism}{单态射}、为\Term{open-immersion}{开浸入}、为\Term{closed-immersion}{闭浸入}、为\Term{quasi-compact}{拟紧}\Term{immersion}{浸入}、\Term{affine}{仿射}、\Term{quasi-affine}{拟仿射}、\Term{finite}{有限}、\Term{quasi-finite}{拟有限}、整、\Term{flat}{平坦}、\Term{faithfully-flat}{忠实平坦}、\Term{net}{净}、\Term{smooth}{光滑}、\Term{etale}{平展}

（EGA IV 2.6.4、2.7.1 及 EGA IV 17.7.1）。另一方面，在 $(\mathrm{Sch})$ 上配备\Term{pretopology}{预拓扑}，使对每个\Term{scheme}{概形} $X\in U$，$\operatorname{Cov}(X)$ 由\Term{surjection}{满射}的\Term{morphism}{态射}族 $(f_i:X_i\to X)_{i\in I}$ 构成，且各 $f_i$ \Term{flat}{平坦}并\Term{locally-finite-presentation}{局部有限表现}。则以下每个性质在上方都具有\Term{local-nature}{局部性质}：

\Term{locally-finite-type}{局部有限型}、\Term{locally-finite-presentation}{局部有限表现}、\Term{finite-type}{有限型}、\Term{flat}{平坦}、\Term{net}{净}、\Term{smooth}{光滑}、\Term{etale}{平展}

（EGA IV 17.7.5、17.7.7）。

\begin{SGAresult}{习题}{8.6}
设 $f:E\to F$ 为 $U$-\Term{topos-morphism}{拓扑斯态射}。考虑关于自变量 $X\in\operatorname{ob}F$ 的关系：诱导\Term{morphism}{态射} $E_{/f^*(X)}\to F_{/X}$ 为\Term{topos}{拓扑斯}\Term{equivalence}{等价}。证明，此关系具有\Term{local-nature}{局部性质}。
\end{SGAresult}

\begin{SGAresult}{习题}{8.7}
（\Term{topos}{拓扑斯}的\Term{partition}{分割}，\Term{topos}{拓扑斯}的和。）设 $E$ 为 $U$-\Term{topos}{拓扑斯}。若\Term{canonical}{典范}\Term{morphism}{态射} $\coprod_{i\in I}e_i\to e$ 为\Term{isomorphism}{同构}，即（第 II 讲 4.6.2）$e$ 为各 $e_i$ 的\Term{supremum}{上确界}，且 $i\ne j$ 蕴含 $e_i\cap e_j=\varnothing_E$，则称 $E$ 的\Term{open-subset}{开集}族 $(e_i)_{i\in I}$，即 $E$ 的\Term{final-object}{终对象} $e$ 的\Term{subobject}{子对象}族，为 $e$（或\Term{topos}{拓扑斯} $E$）的一个\Term{partition}{分割}。

 a) 证明，$E$ 的\Term{object}{对象}族 $(e_i)_{i\in I}$ 为 $E$ 的\Term{partition}{分割}的充要条件是，由\Term{localization-functor}{局部化函子} $E\to E_{/e_i}$ 定义的\Term{functor}{函子}
\[E\longrightarrow\prod_i E_{/e_i}\tag{8.7.1}\]
为\Term{equivalence-of-categories}{范畴等价}。

 b) 反过来，设 $(E_i)_{i\in I}$ 为一族 $U$-\Term{topos}{拓扑斯}，$I\in U$。证明，$E=\prod_{i\in I}E_i$ 为 $U$-\Term{topos}{拓扑斯}（称为\Term{topos}{拓扑斯}族 $(E_i)_{i\in I}$ 的和\Term{topos}{拓扑斯}）。定义 $E$ 的一个\Term{partition}{分割} $(e_i)_{i\in I}$ 及\Term{equivalence-of-categories}{范畴等价} $E_{/e_i}\approx E_i$，使\Term{projection}{投影}\Term{functor}{函子} $E\to E_i$ 等同于\Term{localization-functor}{局部化函子} $E\to E_{/e_i}$。

\SourcePage{258}
 c) 设 $E$ 为\Term{topos}{拓扑斯}族 $(E_i)_{i\in I}$ 的和\Term{topos}{拓扑斯}。证明，对每个 $i\in I$，\Term{projection}{投影}\Term{functor}{函子} $E\to E_i$ 形如 $u_i^*$，其中
\[u_i:E_i\longrightarrow E\tag{8.7.2}\]
为\Term{topos-morphism}{拓扑斯态射}。证明，对任意 $U$-\Term{topos}{拓扑斯} $F$，\Term{functor}{函子} $v\mapsto(v\circ u_i)_{i\in I}$
\[\mathscr{H}om\mathrm{top}(E,F)\longrightarrow\prod_i\mathscr{H}om\mathrm{top}(E_i,F)\tag{8.7.3}\]
为\Term{equivalence-of-categories}{范畴等价}。

 d) 设 $(e_i)_{i\in I}$ 为\Term{topos}{拓扑斯} $E$ 的\Term{partition}{分割}。对任意\Term{topos-morphism}{拓扑斯态射} $g:F\to E$，令 $f_i=g^*(e_i)$，则族 $(f_i)_{i\in I}$ 为 $F$ 的\Term{partition}{分割}，诱导\Term{morphism}{态射} $g_i:E_{/e_i}\to F_{/f_i}$ 可用来在唯一\Term{isomorphism}{同构}意义下重构 $g$；考虑 (8.7.1) 型的\Term{equivalence}{等价}后，\Term{functor}{函子} $g^*$ 等同于\Term{functor}{函子} $g_i^*$ 的\Term{cartesian-product}{笛卡儿积}。由此用形如 $\mathscr{H}om\mathrm{top}(F',E_i)$ 的\Term{category}{范畴}给出\Term{category}{范畴} $\mathscr{H}om\mathrm{top}(F,E)$ 的完整描述，其中 $F'$ 为形如 $F_{/f}$ 的\Term{topos}{拓扑斯}（$f$ 为 $F$ 的\Term{final-object}{终对象}的\Term{direct-summand}{直和项}），而 $E_i=E_{/e_i}$。

 e) 特别地，若 $F$ \Term{connected-nonempty}{连通非空}（见 4.3.5），即对于 $F$ 的任意\Term{partition}{分割} $(f_i)_{i\in I}$，恰有一个 $i\in I$ 使 $f_i\ne\varnothing_F$，证明由\Term{topos-morphism}{拓扑斯态射} (8.7.2) \Term{derivation}{导出}的\Term{canonical}{典范}\Term{functor}{函子}
\[\coprod_{i\in I}\mathscr{H}om\mathrm{top}(F,E_i)\longrightarrow\mathscr{H}om\mathrm{top}(E)\tag{8.7.4}\]
为\Term{equivalence-of-categories}{范畴等价}。更特别地，\Term{topos-morphism}{拓扑斯态射}族 (8.7.2) 诱导\Term{equivalence-of-categories}{范畴等价}
\[\coprod_{i\in I}\operatorname{Point}(E_i)\xrightarrow{\approx}\operatorname{Point}(E).\]

 f) $E$ 的\Term{open-subset}{开集}之集的一个\Term{subset}{子集} $I$，若以 $I$ 为指标的恒等族 $(e_i)_{i\in I}$ 为\Term{partition}{分割}，且各 $e_i\ne\varnothing_E$，则称为 $E$ 的\Term{reduced}{约化}\Term{partition}{分割}。证明，\Term{reduced}{约化}\Term{partition}{分割}之间的细化关系（第 I 讲 4.3.2）使\Term{reduced}{约化}\Term{partition}{分割}之集 $P$ 成为一个 $U$-小的\Term{ordered-set}{有序集}，且其中任意两个元素有\Term{supremum}{上确界}。由此得到严格\Term{projective-system}{投射系} $(I_p)_{p\in P}$，将其看作\Term{pro-set}{pro-集合}并记为 $\pi_0(E)$。称它为\Term{topos}{拓扑斯} $E$ 的前 $\pi_0$。

 g) 证明，$\pi_0(E)$ \Term{pro-representable-functor}{pro-表出}与 $E$ 到\Term{punctual-topos}{点拓扑斯} $P$ 的\Term{canonical}{典范}\Term{morphism}{态射} $f:E\to P$（4.3）\Term{correspondence}{对应}的\Term{functor}{函子}
\[T\longmapsto f_*f^*(T)=\Gamma(E,T_E):(U\text{-}\mathrm{Ens})\longrightarrow(U\text{-}\mathrm{Ens}).\]
特别地，此\Term{functor}{函子}\Term{representable}{可表}的充要条件是 $\pi_0(E)$ \Term{essentially-constant-ind-object}{本质常值}，即（作为\Term{pro-set}{pro-集合}）\Term{isomorphism}{同构}于一个“普通”\Term{set}{集合}，仍将此\Term{set}{集合}记为 $\pi_0(E)$。$\pi_0(E)$ 只有一个点的充要条件
\SourcePage{259}
是 $E$“\Term{connected-nonempty}{连通非空}”。$\pi_0(E)$ 为空的充要条件是 $E$ 为“\Term{empty-topos}{空拓扑斯}”（2.2）。

 h) 假定 $E$ \Term{quasi-compact}{拟紧}，即其\Term{final-object}{终对象} $e$ 的每个\Term{covering}{覆盖}都有\Term{finite}{有限}子\Term{covering}{覆盖}。证明，此时 $\pi_0(E)$ 为\Term{profinite-set}{投有限集}，因而（借助\Term{finite}{有限}集\Term{category}{范畴}的\Term{pro-object}{pro-对象}与紧致全不\Term{connected}{连通}\Term{space}{空间}之间熟知的\Term{equivalence-of-categories}{范畴等价}）可以等同于一个紧致全不\Term{connected}{连通}\Term{space}{空间}，仍记为 $\pi_0(E)$。

 i) 设 $X$ 为\Term{topological-space}{拓扑空间}，$\pi_0(X)$ 为 $X$ 的\Term{connected-component}{连通分支}之集，看成一个常\Term{pro-set}{pro-集合}。定义\Term{pro-set}{pro-集合}的\Term{canonical}{典范}\Term{morphism}{态射}
\[\pi_0(X)\longrightarrow\pi_0(\operatorname{Top}(X)),\tag{8.7.5}\]
或\Term{equivalence}{等价}地，定义\Term{canonical}{典范}映射
\[\pi_0(X)\longrightarrow\varprojlim\pi_0(\operatorname{Top}(X)).\tag{8.7.6}\]
证明，第一个\Term{morphism}{态射}为\Term{isomorphism}{同构}，当且仅当 $X$ \Term{homeomorphism}{同胚}于\Term{connected}{连通}\Term{space}{空间}的和\Term{space}{空间}（特别地，当 $X$ \Term{locally-connected}{局部连通}时即如此）。

 j) 假定 $X$ 为\Term{quasi-compact}{拟紧}\Term{scheme}{概形}。证明此时 (8.7.6) 为\Term{bijection}{双射}。

 k) 建立\Term{morphism}{态射} (8.7.5) 和 (8.7.6) 关于 $X$ 的“\Term{functor}{函子}性”，并先明确说明此表述的意义。

 l)（比较 7.6 i)。）证明关于 $U$-\Term{topos}{拓扑斯} $E$ 的以下两个条件\Term{equivalence}{等价}：(i) 从 $E$ 到\Term{punctual-topos}{点拓扑斯}的\Term{canonical}{典范}\Term{morphism}{态射}为“本质”的（习题 7.6 a)），即从 $(U\text{-}\mathrm{Ens})$ 到 $E$ 的\Term{functor}{函子} $I\mapsto I_E$ 与以属于 $U$ 的\Term{set}{集合}为指标的积交换；(ii) 对 $E$ 的任意\Term{object}{对象} $X$，\Term{induced-topos}{诱导拓扑斯} $E_{/X}$ 的 $\pi_0(E_{/X})$ 为普通\Term{set}{集合}，即 $X$ \Term{isomorphism}{同构}于一族 $E$ 的\Term{connected}{连通}\Term{object}{对象}的和。此时称 $E$ 为\Term{locally-connected-topos}{局部连通拓扑斯}（比较第 IV 讲 2.7.5）。
\end{SGAresult}

\begin{SGAresult}{习题}{8.8}
（\Term{dominant-morphism}{支配}\Term{topos-morphism}{拓扑斯态射}。）

 a) 设 $f:E'\to E$ 为\Term{topos-morphism}{拓扑斯态射}。证明以下条件\Term{equivalence}{等价}：

 (i) $f_*(\varnothing_{E'})=\varnothing_E$（其中 $\varnothing_E,\varnothing_{E'}$ 表示 $E,E'$ 的\Term{initial-object}{始对象}）。

 (ii) 对 $E$ 的每个“\Term{nonempty}{非空}”\Term{object}{对象} $X$，即不\Term{isomorphism}{同构}于 $\varnothing_E$ 的\Term{object}{对象}，$f^*(X)$ 为 $E'$ 的\Term{nonempty}{非空}\Term{object}{对象}。

 (ii') 同 (ii)，其中 $X$ 为 $E$ 的\Term{final-object}{终对象} $e_E$ 的\Term{subobject}{子对象}，即\Term{topos}{拓扑斯} $E$ 的\Term{open-subset}{开集}。

此时称\Term{topos-morphism}{拓扑斯态射} $f$ 为\Term{dominant-morphism}{支配}的。

\SourcePage{260}
 b) 设 $f:X'\to X$ 为\Term{topological-space}{拓扑空间}的\Term{continuous-map}{连续映射}，从而有\Term{topos-morphism}{拓扑斯态射} $\operatorname{Top}(f):\operatorname{Top}(X')\to\operatorname{Top}(X)$（4.1）。证明，$\operatorname{Top}(f)$ \Term{dominant-morphism}{支配}当且仅当 $f$ \Term{dominant-morphism}{支配}，即 $f(X')$ 在 $X$ 中\Term{dense}{稠密}。

 c) 证明，若 $f:E'\to E$ 为“\Term{conservative-functor}{保守}”的\Term{topos-morphism}{拓扑斯态射}（6.4.0），即 $f^*$ \Term{faithful-functor}{忠实}，则 $f$ \Term{dominant-morphism}{支配}。证明逆命题未必成立，甚至对于\Term{localization-morphism}{局部化态射} $E_{/U}\to E$（$U$ 为\Term{topos}{拓扑斯} $E$ 的\Term{open-subset}{开集}）也未必成立。（证明此时 $f$ \Term{conservative-functor}{保守}仅当 $U$ 为 $E$ 的\Term{final-object}{终对象}（因而 $f$ 为\Term{topos}{拓扑斯}\Term{equivalence}{等价}），并利用例 b)。）

 d) 设 $f:X'\to X$ 为\Term{topos}{拓扑斯} $E$ 中的\Term{arrow}{箭头}。若\Term{correspondence}{对应}的\Term{induced-topos}{诱导拓扑斯}\Term{morphism}{态射} $E_{/X'}\to E_{/X}$（5.5.2）\Term{dominant-morphism}{支配}，则称 $f$ 为\Term{topos}{拓扑斯} $E$ 中的\Term{dominant-morphism}{支配}\Term{morphism}{态射}。证明，此性质只依赖于 $X'$ 在 $X$ 中的像 $f(X')$，即只依赖于它作为 $X$ 的\Term{subobject}{子对象}的\Term{ordered-set}{有序集} $\operatorname{ouv}(X)$ 的一个元素。证明，\Term{morphism}{态射} $E_{/X'}\to E_{/X}$ \Term{conservative-functor}{保守}当且仅当 $f$ 为\Term{covering}{覆盖}\Term{morphism}{态射}。
\end{SGAresult}
